% The manuscript (status: draft). Paths such as code/fourier/existence.py and data/fourier-stability-N8.json are
% relative to the paper's folder; the canonical copies are in research/cardiac-cycle-certificates/ of GENChase.
\documentclass[11pt]{article}
\usepackage[letterpaper,margin=1in]{geometry}
\usepackage[T1]{fontenc}
\usepackage{lmodern}
\usepackage{amsmath,amssymb,amsthm}
\usepackage{array}
\usepackage{booktabs}
\usepackage{microtype}
\setlength{\emergencystretch}{2em}
\usepackage{xurl}
\usepackage[hidelinks]{hyperref}

\hypersetup{
  pdftitle={Stable Rotating Waves in Rings of a Modified Ventricular Myocyte Model Near a Hopf Point: Computer-Assisted Proofs in Fourier Space},
  pdfauthor={Chase Hendrick}}

\newtheorem{theorem}{Theorem}
\renewcommand{\thetheorem}{\Alph{theorem}}
\newtheorem{lemma}{Lemma}[section]
\newtheorem{proposition}[lemma]{Proposition}
\newtheorem{thmx}[lemma]{Theorem}
\newtheorem{corollary}[lemma]{Corollary}
\theoremstyle{definition}
\newtheorem{definition}[lemma]{Definition}
\newtheorem{remark}[lemma]{Remark}

\newcommand{\R}{\mathbb{R}}
\newcommand{\C}{\mathbb{C}}
\newcommand{\Z}{\mathbb{Z}}
\newcommand{\re}{\operatorname{Re}}
\newcommand{\im}{\operatorname{Im}}
\newcommand{\spec}{\operatorname{spec}}
\newcommand{\Ran}{\operatorname{Ran}}
\newcommand{\rank}{\operatorname{rank}}
\newcommand{\one}{\mathbf{1}}
\newcommand{\Dsp}{\mathcal{D}}
\newcommand{\Xsp}{\mathcal{X}}
\newcommand{\Lop}{\mathcal{L}}
\DeclareUrlCommand\file{\urlstyle{tt}}

\title{\Large\bfseries Stable Rotating Waves in Rings of a Modified Ventricular Myocyte Model Near a Hopf Point:\\ Computer-Assisted Proofs in Fourier Space}
\author{Chase Hendrick\\[0.2em] {\small Independent Researcher}\\ {\small\href{mailto:chase@hendrickresearch.com}{\texttt{chase@hendrickresearch.com}}}\\ {\small ORCID \href{https://orcid.org/0009-0002-9754-6087}{0009-0002-9754-6087}}}
\date{Draft of October 1, 2026}

\begin{document}
\maketitle

\begin{abstract}
We study an 18-dimensional modification, due to Erhardt, of the ten Tusscher--Panfilov 2006 endocardial ventricular cell model, with reduced repolarization reserve and the slow delayed rectifier conductance $G_{Ks} = 0.0275$~nS/pF, a value about 1.5\% below the supercritical Hopf point that Erhardt computed numerically for this model. For the single cell we give two computer-assisted proofs, sharing no library, that a periodic orbit exists and is locally orbitally asymptotically stable: a time-domain proof with CAPD, which encloses the period in $[53.5855190480139, 53.585519630722438]$~ms and bounds the 17 nontrivial Floquet multipliers of its orbit in modulus by $0.998642$, and a space-time Fourier proof in Arb ball arithmetic, which encloses the period of its orbit in an interval of width less than $2 \times 10^{-25}$~ms and bounds its 17 nontrivial multipliers by $0.99785888$. An exact rational check shows that a point of the second orbit lies in the ball that the first proof shows the return map to map into itself as a contraction, so the two proofs are about the same orbit. For rings of $N = 8$, $16$, $32$ and $64$ identical cells with voltage-only diffusive coupling of strength $N^2/64000$ per ms we prove, in Fourier space, that a rotating 1-wave $x_j(t) = \phi(\omega t + 2\pi j/N)$ exists, is locally unique, has minimal period enclosed in an interval of width less than $2 \times 10^{-25}$~ms, is not synchronous, and is locally exponentially orbitally stable with asymptotic phase: the Floquet multiplier $1$ is algebraically simple, and the other $18N - 1$ multipliers have modulus less than $e^{-\delta T}$ with $\delta = 5 \times 10^{-6}$ per ms. Existence is proved by a radii-polynomial argument in a weighted $\ell^1$ space, with rigorous strip covers, aliasing bounds and a polydisc Cauchy majorant for the non-polynomial ionic currents. Stability is proved through one Hill operator: we show that its spectrum on a half-open strip of height $\omega N$ gives every Floquet multiplier of the ring with its algebraic multiplicity, and we exclude spectrum from $\{\re \mu \ge -\delta\}$ apart from a simple eigenvalue $0$ by a Riesz-projection homotopy with a Schur-complement small-gain test and an explicit tail resolvent bound. The identification of these orbits with the Hopf branch is numerical, and nothing is claimed for a continuum cable or for tissue.

\medskip
\noindent\textbf{Keywords:} cardiac cell model; modified ten Tusscher--Panfilov model; ring of coupled oscillators; rotating wave; Floquet multipliers; Hill operator; computer-assisted proof; ball arithmetic; radii polynomial.\\
\textbf{MSC 2020:} 37C27, 34C25, 65G20, 92C30, 37M20.
\end{abstract}

\section{Introduction}\label{sec:intro}

Detailed ionic models of cardiac cells are systems of a few dozen nonlinear ordinary differential equations whose right-hand sides contain exponentials, logarithms (Nernst potentials) and square roots. Their bifurcations are studied by numerical continuation. Erhardt and Solem analysed reduced versions of the ten Tusscher--Panfilov 2006 (TP06) model \cite{tp2006} in this way \cite{es2021,es2022}, K\"ugler and coauthors related early afterdepolarizations to period doubling cascades and to folded nodes \cite{keb2017,keb2018}, and Erhardt \cite{erhardt2025} studied an 18-dimensional modification of the TP06 endocardial cell with reduced repolarization reserve, continued in the slow delayed rectifier conductance $G_{Ks}$. For that model he reports a supercritical Hopf bifurcation at $G_{Ks} = 0.027907858929580$~nS/pF with stable bifurcating cycles \cite[Table~2]{erhardt2025}.

This paper gives computer-assisted proofs for that model at $G_{Ks} = 0.0275$, about $1.46\%$ below the Hopf value, for the single cell and for rings of $N$ identical cells coupled diffusively through the membrane potential. The ring is the simplest spatial object built from these cells: a periodic chain $x_0, \dots, x_{N-1}$ with
\begin{equation}\label{eq:ring}
  \dot x_j = f(x_j) + c\, E\,(x_{j-1} - 2x_j + x_{j+1}), \qquad j \in \Z_N, \qquad c = \frac{N^2}{64000}\ \text{per ms},
\end{equation}
where $f$ is the cell's vector field and $E$ the projection on the voltage. The coupling $c = N^2 D$ with $D = 1/64000$ is that of a ring of unit length discretized into $N$ cells. The ring is equivariant under the cyclic shift, and a rotating wave (a discrete traveling wave) is a solution $x_j(t) = \phi(\omega t + 2\pi j/N)$ with $\phi$ $2\pi$-periodic. That such waves exist near a Hopf point of a ring of identical cells is the generic expectation of $\Z_N$-equivariant Hopf theory \cite{gs1985}, \cite[Ch.~XVIII, Sec.~4]{gss1988}; we make no claim that their existence is surprising. What we add is a proof, at an explicit parameter value and explicit coupling strengths, of existence, local uniqueness, the minimal period and local exponential orbital stability, for the full 18-dimensional nonlinear cells.

\paragraph{Results.} Theorem~\ref{thm:cell} concerns the single cell. It has two parts, proved independently: one in the time domain with CAPD \cite{capd2021}, by a contraction argument for the Poincar\'e map, and one in Fourier space with Arb \cite{johansson2017}, by the method of this paper with $N = 1$. The two proofs share no library, and an exact rational check (Lemma~\ref{lem:link}) shows that they enclose the same orbit. We call the existence proof in Fourier space Stage~E (Section~\ref{sec:existence}) and the stability proof Stage~S (Section~\ref{sec:stability}), as the programs and their records do. Theorem~\ref{thm:ring} concerns the rings with $N = 8, 16, 32, 64$. Its existence part (Section~\ref{sec:existence}) is a radii-polynomial (Newton--Kantorovich) argument for the Fourier coefficients of the profile $\phi$ in a weighted $\ell^1$ space. In Fourier space the ring enters only through a scalar damping $d_m = 4c\sin^2(\pi m/N)$ on the voltage of mode $m$, so the cost does not grow with the dimension $18N$ of the ring: each existence proof takes about two minutes. Its stability part (Section~\ref{sec:stability}) rests on a lemma (Theorem~\ref{thm:hill}) that the Floquet multipliers of the rotating wave, with algebraic multiplicities, are the numbers $e^{\mu T}$ for $\mu$ in the spectrum of a single Hill operator $H_0$ in a half-open strip of height $\omega N$. Spectrum of $H_0$ is then excluded from the half plane $\{\re \mu \ge -\delta\}$, except a simple eigenvalue at $0$ and its translates by $i\omega N\Z$, by counting the rank of a Riesz projection along a homotopy (Theorem~\ref{thm:cert}). The stability proofs take between half a minute and eight minutes.

\paragraph{Related work.} Space-time Fourier existence proofs for periodic orbits by the radii-polynomial approach are due to Gameiro and Lessard \cite{gl2017}, who proved periodic orbits of the Kuramoto--Sivashinsky equation and enclosed individual Floquet exponents by a second contraction on the invariance equation \cite[Sec.~7]{gl2017}. That proves instability; they state that their approach could only prove that some solutions are unstable, and that extending it to prove stability ``is the subject of current research'' \cite[Sec.~1]{gl2017}. The exclusion and counting step of Section~\ref{sec:stability} is how we supply the missing part for our Hill operator. Other routes to the stability of a periodic orbit exist: Castelli and Lessard \cite{cl2013} computed the Floquet normal form $\Phi(t) = Q(t)e^{Rt}$ of periodic orbits of ODEs by rigorous numerics, which gives all Floquet exponents; a time-domain contraction for the Poincar\'e map, as in Section~\ref{sec:capd}, is another; and Church, Dai, H\'enot, Lappicy and Vassena \cite{cdhlv2026} proved global families of stable periodic orbits in a model of competing predators by computer; their abstract does not say how stability is certified. Figueras, Gameiro, Lessard and de la Llave \cite{fgld2017} gave a framework for the numerical computation and a posteriori verification of invariant objects of evolution equations. Church and Lessard \cite{cl2022} verified Hopf bifurcations in functional differential equations of mixed type, the class to which the profile equation of a lattice traveling wave belongs, Arioli and Koch \cite{ak2020} constructed traveling waves of the FPU chain by a computer-assisted proof, and Kuehn and Queirolo \cite{kq2022} proved Hopf bifurcation points and periodic orbits of recurrent neural networks by computer. Bayer and Leine \cite{bl2026} proved an explicit a priori bound on the truncation error of the Koopman--Hill approximation of the monodromy matrix of a linear time-periodic system, for coefficients with $\|J_k\| \le a e^{-b|k|}$ and $b > \ln 2$ \cite[Theorem~4]{bl2026}, and inferred multiplier locations from pseudospectra in floating point; our tail bound is a posteriori and acts on the resolvent of the Hill operator, which lets us count eigenvalues in ball arithmetic. Johnson and Zumbrun \cite{jz2012} proved convergence of Hill's method for nonselfadjoint spatially periodic operators; we cite it as background on Hill truncation only, since it gives no explicit error bounds and does not cover the operator $d/dt - A(t)$ of a time-periodic Floquet problem.

The reduction of a ring's Floquet problem to Fourier blocks combines three classical ingredients: the $\Z_N$-isotypic decomposition of rings of identical cells \cite{gs1985}, \cite[Ch.~XVIII, Sec.~4]{gss1988}; the spatio-temporal symmetry of a discrete wave, which makes the monodromy an $N$-th power of a twisted map, as in Rucklidge and Silber \cite{rs1998} and de Wolff \cite{dewolff2023}; and the discrete Laplacian eigenvalues $4\sin^2(\pi m/N)$, familiar from the master stability function of Pecora and Carroll \cite{pc1998}. We state the resulting identity as a lemma and prove it for completeness. Paullet and Ermentrout \cite{pe1994} studied stable rotating waves in two-dimensional discrete active media analytically, and Ermentrout \cite{ermentrout1992} gave conditions for the stability of phase-locked solutions of weakly coupled oscillator arrays (we read his Theorem~3.1 and Lemma~3.2); Ashwin and Swift \cite{as1992} and Hoppensteadt and Izhikevich \cite{hi1997} treat weakly coupled identical oscillators and phase reduction. For piecewise-linear Chua--Yang rings, Di Marco, Forti, Garay, Koller and Pancioni \cite{dmfgkp2016} obtained, according to their abstract, the Floquet multipliers of an unstable, metastable rotating wave. Reentry on rings of excitable (not self-oscillating) cardiac tissue has a long numerical and analytical history, for example Courtemanche, Glass and Keener \cite{cgk1993} and ten Tusscher and Panfilov \cite{tp2006}; the waves of this paper are of a different kind, phase waves on rings of self-oscillating cells. For the FitzHugh--Nagumo traveling-wave ODE, Czechowski and Zgliczy\'nski \cite{cz2016} proved periodic orbits by computer, without stability, and, according to them, Arioli and Koch \cite{ak2015} proved the existence and stability of the FitzHugh--Nagumo traveling pulse by computer. Kapela and Zgliczy\'nski \cite{kz2003} proved choreographies of the $N$-body problem by computer; a choreography has the same cyclic-shift symmetry $x_{j+1}(t) = x_j(t + T/N)$ as a rotating wave, in a Hamiltonian system without diffusive coupling.

\paragraph{Novelty.} The searches logged in this project's research ledger (Section~\ref{sec:limits}) found no earlier computer-assisted proof of a periodic orbit of a detailed ionic cardiac cell model, and no earlier computer-assisted proof of a rotating wave in a ring of diffusively coupled cells. We make no claim beyond that sentence; the nearest computer-assisted precedents for the symmetry and for stability are cited above. The oscillation of the single cell is not new: Erhardt's continuation \cite{erhardt2025} predicts it.

\paragraph{Labels.} Theorems~\ref{thm:cell} and~\ref{thm:ring} are computer-assisted: every inequality they rest on is decided in interval or ball arithmetic by the programs of Section~\ref{sec:reproduce}, and every step from those inequalities to the statements is proved in this paper. The lemmas are proved here; three standard facts about closed operators with compact resolvent are cited from Kato's book (Section~\ref{sec:hill}). Section~\ref{sec:numerics} collects numerical observations, labelled as such; none of them is used in a proof or stated as a theorem. The checks made of this work are listed in Section~\ref{sec:limits}.

\section{The model and the ring}\label{sec:model}

\subsection{The cell}\label{sec:cell}

The vector field $f: U \subset \R^{18} \to \R^{18}$ is the function \texttt{fun\_eval} of the file \texttt{bifurcation\ analysis/TP06\_18d\_endo\_bif.m} in Erhardt's MIT-licensed repository \file{andreerhardt/cardiac-dynamics-of-a-human-ventricular-tissue-model-with-focus-on-early-afterdepolarizations}, commit \texttt{dc78f86fd218418e029ec43d945bcd0fc54b9f1e}, file SHA-256 \texttt{a50f6c08b4360dd257cce389a39ae72fda51e3642641bf5b8e5fced6c2225670}, which accompanies \cite{erhardt2025}. The state is $(V, X_{r1}, X_{r2}, X_s, m, h, j, d, f, f_2, f_{Cass}, s, r, R', \mathrm{Ca}_i, \mathrm{Ca}_{sr}, \mathrm{Ca}_{ss}, \mathrm{Na}_i)$, in mV and mM, and time is in ms. It is Erhardt's 18-state modification of the TP06 endocardial cell \cite{tp2006}, and we call it that rather than ``the TP06 cell'', because it differs from the published TP06 cell in four ways:
\begin{enumerate}
\item the intracellular potassium concentration $K_i$ is a parameter held at $138.3$~mM, so 19 states become 18;
\item the Heaviside switch at $V = -40$~mV in the rates of the gates $h$ and $j$ is replaced by $u(V) = 1/(1 + e^{-5(V + 40)})$, which makes $f$ real analytic;
\item Erhardt's reduced-repolarization parameters: $G_{Kr} = 0.0153$ ($0.1$ times the endocardial value), $G_{CaL} = 0.000199$ ($5$ times), with $G_{Ks}$ the continuation parameter, here $G_{Ks} = 0.0275$; also $G_{Na} = 14.838$, $G_{K1} = 5.405$ and no stimulus;
\item Erhardt's capacitance convention: the parameter $C_m = 1$ divides $dV/dt$ and also multiplies the fluxes of $\mathrm{Ca}_i$, $\mathrm{Ca}_{ss}$ and $\mathrm{Na}_i$, where the original authors' convention (and the CellML version) uses the numerical value $0.185$ for the cell capacitance (185~pF); so every concentration flux is $1/0.185 \approx 5.405$ times its value in that convention. We state only this ratio and make no claim about the physical units of Erhardt's $C_m$.
\end{enumerate}
The domain $U$ is the open set on which every logarithm argument is positive, every square-root argument is positive and every divisor is nonzero. In the proofs every decimal constant of the source enters as an exact rational (Section~\ref{sec:reproduce}), so the theorems are about the model as written, not about its floating-point rounding. All computations use the scaled variables $z = \Sigma^{-1} x$, $\Sigma = \operatorname{diag}(2^{e_1}, \dots, 2^{e_{18}})$ with the scale exponents $(-2, 0, -5, -1, 0, -28, -28, 0, -4, -2, -1, -8, -5, -1, -10, 2, -3, 3)$ (\file{code/model/scales.txt}); a constant diagonal change of variables commutes with $E$ and changes no eigenvalue or multiplier, and we keep writing $f$ for the field in scaled variables where no confusion arises.

\subsection{The ring and its symmetry}\label{sec:ringsym}

For $N \ge 1$ the ring is \eqref{eq:ring} on $(\R^{18})^N$, with $E = e_V e_V^{\mathsf T}$ the projection on the voltage component; for $N = 1$ the coupling vanishes and \eqref{eq:ring} is the cell. Write $F$ for its vector field and $Q$ for the shift $(Qx)_j = x_{j+1}$. Then $F(Qx) = QF(x)$ and $Q^N = I$.

Let $\omega > 0$, $T = 2\pi/\omega$, $\tau = T/N$, and let $\phi: \R \to \R^{18}$ be $2\pi$-periodic, real analytic and nonconstant, with $\phi(\theta) = \sum_{m \in \Z} a_m e^{im\theta}$. Write $\theta_j(t) = \omega t + 2\pi j/N$ and $(\Delta_N p)(\theta) = p(\theta - 2\pi/N) - 2p(\theta) + p(\theta + 2\pi/N)$. If $\phi$ solves the profile equation
\begin{equation}\label{eq:profile}
  \omega\,\phi'(\theta) = f(\phi(\theta)) + c\,E\,(\Delta_N\phi)(\theta) \qquad (\theta \in \R),
\end{equation}
then $x_j(t) = \phi(\theta_j(t))$ solves \eqref{eq:ring}, because $x_{j\pm1}(t) = \phi(\theta_j(t) \pm 2\pi/N)$, and indices modulo $N$ are consistent because $\phi$ is $2\pi$-periodic. Moreover $x(t + \tau) = Qx(t)$, since $\theta_j(t + \tau) = \theta_{j+1}(t)$. We call $x$ a rotating 1-wave. On mode $m$ the operator $c\,E\,\Delta_N$ acts as $-d_m E$ with
\begin{equation}\label{eq:damping}
  d_m = 4c\sin^2(\pi m/N), \qquad 0 \le d_m \le 4c, \quad d_{-m} = d_m = d_{m+N} = d_{N-m},
\end{equation}
so in Fourier space \eqref{eq:profile} reads $i\omega m a_m = [f \circ \phi]_m - d_m E a_m$ for every $m \in \Z$. The ring is felt only through the scalars $d_m$.

\section{Results}\label{sec:results}

Let $s$ be the double-precision number nearest to $0.2$, and call the section $\{V = s\}$ (in mV).

\begin{theorem}[the cell; computer-assisted]\label{thm:cell}
Let $N = 1$ and $G_{Ks} = 0.0275$.
\begin{enumerate}
\item[(i)] (CAPD, time domain.) The first-return map $g$ to the section $\{V = s\}$, crossed upward, has a unique fixed point $x_c$ in the explicit ball $B$ of Section~\ref{sec:capd}. The periodic orbit through $x_c$ has period in $[53.5855190480139, 53.585519630722438]$~ms (exact bounds \texttt{0x1.acaf249c533cfp+5} and \texttt{0x1.acaf24ea88f37p+5}), every eigenvalue of $Dg(x_c)$, that is every one of the 17 nontrivial Floquet multipliers, has modulus at most $0.99864150686405151 < 0.998642$, and the orbit is locally orbitally asymptotically stable.
\item[(ii)] (Arb, Fourier space.) There are $\omega_* > 0$ and a real analytic $2\pi$-periodic $\phi_*$ with $\phi_{*,V}(0) = s$ such that $x(t) = \phi_*(\omega_* t)$ is a periodic orbit of minimal period $T = 2\pi/\omega_*$ with
\[ T \in [53.58551933936116920991898,\ 53.58551933936116920991899]~\text{ms} \]
(23-decimal outward roundings of binary ends $1.56 \times 10^{-25}$~ms apart);
the pair $(\omega_*, (a_m))$ is locally unique in the sense of Theorem~\ref{thm:ring}(a) with $N = 1$. The Floquet multiplier $1$ is algebraically simple and the other 17 multipliers have modulus less than $e^{-\delta T} \le 0.9978588747123748493822038$, where $\delta = 23058430092137/2^{59} \approx 4.0000000000000104916 \times 10^{-5}$ per ms; the orbit is locally exponentially orbitally stable with asymptotic phase.
\item[(iii)] (Link.) The orbits of (i) and (ii) are the same: $\Sigma^{-1}\phi_*(0) \in B$, so $x_c$ lies on the orbit of (ii) (Lemma~\ref{lem:link}). The period interval of (ii) lies inside that of (i).
\end{enumerate}
\end{theorem}

Each of (i) and (ii) assumes that its program evaluates Erhardt's function with exact decimals (the CAPD field in (i), the Arb field in (ii)); (iii) uses both assumptions together, so a fault in either translation would also break the link. The two programs were compared at 104 points by the tests (Section~\ref{sec:reproduce}), not proved equal (Section~\ref{sec:limits}).

\begin{theorem}[rings; computer-assisted]\label{thm:ring}
Let $G_{Ks} = 0.0275$, $N \in \{8, 16, 32, 64\}$ and $c = N^2/64000$ per ms.
\begin{enumerate}
\item[(a)] (Existence and local uniqueness.) There are $\omega_* > 0$ and a real $2\pi$-periodic $\phi_*$, analytic on the strip $|\im\theta| < 1/4$, with $\phi_{*,V}(0) = s$, such that $x_j(t) = \phi_*(\omega_* t + 2\pi j/N)$ solves \eqref{eq:ring}. The pair $(\omega_*, (a_m)_{m\in\Z})$ formed by $\omega_*$ and the Fourier coefficients of $\Sigma^{-1}\phi_*$ (the profile in the scaled variables of Section~\ref{sec:cell}) lies within $r_{\mathrm{ex}}$ of an explicit centre in the norm $\max(|\omega|, \max_k \sum_m |a_{k,m}| e^{|m|/4})$, and it is the only zero of the map \eqref{eq:F} below in the ball of radius $r_{\mathrm{un}}$ about that centre, where $r_{\mathrm{un}}$ is the largest double not above $10^{-12}$ ($r_{\mathrm{un}} \approx 9.9999999999999997989 \times 10^{-13}$); $r_{\mathrm{ex}}$ is given in Table~\ref{tab:existence}.
\item[(b)] (Period, minimal period, wave number.) $T = 2\pi/\omega_*$ lies in a certified interval of width less than $2 \times 10^{-25}$~ms (between $1.49 \times 10^{-25}$ and $1.57 \times 10^{-25}$ for the five values of $N$), which Table~\ref{tab:existence} prints rounded outward to 23 decimals. Each cell has minimal period $T$, the solution is not synchronous (it is not the case that $x_j = x_0$ for all $j$), and it is a rotating 1-wave: $x_j(t) = x_0(t + jT/N)$, and no relation $x_j(t) = x_0(t + jkT/N)$ for all $j$ holds with $k \not\equiv 1 \pmod N$.
\item[(c)] (Stability.) Let $\delta = 5764607523035/2^{60} \approx 5.0000000000006636358 \times 10^{-6}$ per ms. The monodromy matrix of the linearization along the wave over the period $T$ has the eigenvalue $1$ algebraically simple, and its other $18N - 1$ eigenvalues (the nontrivial Floquet multipliers, with algebraic multiplicity) have modulus less than $e^{-\delta T}$, which is at most the bound of Table~\ref{tab:stability}. The shift-reduced map $Q^{-1}\varphi_{\tau}$, $\tau = T/N$, has derivative at $x(0)$ with eigenvalue $1$ algebraically simple and all other eigenvalues of modulus less than $e^{-\delta\tau}$. The orbit is locally exponentially orbitally stable with asymptotic phase: for every $\delta' < \delta$ there are $\epsilon_0 > 0$ and $C$ such that every solution starting within distance $\epsilon_0$ of the orbit $O$ satisfies, for some $\sigma \in \R$ and all $t \ge 0$, $|\varphi_t(x^0) - x(t + \sigma)| \le C e^{-\delta' t} \operatorname{dist}(x^0, O)$.
\end{enumerate}
\end{theorem}

Here $\varphi_t$ is the flow of \eqref{eq:ring}. The constants $\epsilon_0$ and $C$ are not computed.

\begin{table}[ht]
\centering\small
\begin{tabular}{@{}rlll@{}}
\toprule
$N$ & lower end of the enclosure of $T$ (ms) & $r_{\mathrm{ex}}$ & $Z_1$ \\
\midrule
1 & 53.58551933936116920991898 & $1.7073 \times 10^{-28}$ & 0.20366 \\
8 & 53.58797097681944967415082 & $1.6440 \times 10^{-28}$ & 0.20339 \\
16 & 53.58806910359016923593730 & $1.6415 \times 10^{-28}$ & 0.20338 \\
32 & 53.58809413031803250655154 & $1.6409 \times 10^{-28}$ & 0.20337 \\
64 & 53.58810041831757754104213 & $1.6408 \times 10^{-28}$ & 0.20337 \\
\bottomrule
\end{tabular}
\caption{Existence (Section~\ref{sec:existence}), from the existence records in \texttt{data/}. The printed ends are the exact binary ends stored in the record, about $1.5 \times 10^{-25}$~ms apart, rounded outward to 23 decimals; the upper end is one unit larger in the last printed digit. $r_{\mathrm{ex}}$ (rounded up here) and $Z_1$ (rounded) are the radius at which the radii polynomial is negative and the bound of Lemma~\ref{lem:rp}; the uniqueness radius is $r_{\mathrm{un}} \approx 10^{-12}$ for every $N$, set by the validity radius $r_*$ of the bound $Z_2$.}
\label{tab:existence}
\end{table}

\begin{table}[ht]
\centering\small
\begin{tabular}{@{}rlllll@{}}
\toprule
$N$ & bound on $e^{-\delta T}$ (full period) & bound on $e^{-\delta\tau}$ & $K_e$ & $\theta_T$ & worst (SC) ratio \\
\midrule
1 & 0.9978588747123748493822038 & (same) & 16 & 0.2948 & 0.104 \\
8 & 0.9997320960377930203747358 & 0.9999665080790063923613739 & 20 & 0.2817 & 0.090 \\
16 & 0.9997320955472906096505621 & 0.9999832538686231493032797 & 24 & 0.2519 & 0.084 \\
32 & 0.9997320954221904942011761 & 0.9999916268953467731405854 & 32 & 0.2895 & 0.154 \\
64 & 0.9997320953907589193958104 & 0.9999958134384184917812441 & 48 & 0.2950 & 0.169 \\
\bottomrule
\end{tabular}
\caption{Stability (Section~\ref{sec:stability}), from the stability records in \texttt{data/}: the certified upper bounds, rounded upward from Arb, with $\delta = 23058430092137/2^{59}$ for $N = 1$ and $\delta = 5764607523035/2^{60}$ for the rings; the window half-width $K_e$; the tail constant $\theta_T < 1$ of Lemma~\ref{lem:smallgain}; and the largest ratio of the left side to the right side of the window inequality (SC) over all window columns, which must be below $1$.}
\label{tab:stability}
\end{table}

\begin{remark}\label{rem:same}
The CAPD record \file{data/cell-gks0.0275.json} and the Fourier record \file{data/fourier-existence-N1.json} store both period enclosures, and the second lies inside the first; Lemma~\ref{lem:link} shows that the two orbits are the same. An independent computation in the same project, by a separate pipeline with CAPD on the owner's machine (\file{docs/CARDIAC-HANDOFF-2026-09-30.md} of GENChase; not a publication), enclosed the periods of the cell in $[53.58551856, 53.58552012]$~ms, of the 8-cell wave in $[53.58795907, 53.58798288]$~ms and of the 16-cell wave in $[53.58805554, 53.58808267]$~ms, with all nontrivial multipliers of modulus below $999/1000$ (cell) and $9999/10000$ (rings). The periods of Theorem~\ref{thm:cell}(ii) and Theorem~\ref{thm:ring} for $N = 1, 8, 16$ lie inside these intervals. We use these intervals as a consistency check only.
\end{remark}

\section{Existence in Fourier space}\label{sec:existence}

\subsection{Rigorous Fourier coefficients of a composition}\label{sec:fe}

Let $p(\theta) = \sum_{|m| \le K} a_m e^{im\theta}$ with $a_m \in \C^n$ given as complex balls; every statement below holds for every trigonometric polynomial with coefficients in the balls. Let $f: U \to \C^q$ be holomorphic on an open $U \subset \C^n$, given by an inclusion function $\mathsf F$: a program that maps a list of complex balls $X$ to balls $Y$ or raises an error, such that whenever all of $Y$ is finite, $X \subset U$ and $f(z) \in Y$ for every $z \in X$. A program written with $+, -, \times$, integer powers, $\exp$, division, and logarithms and square roots that raise unless their argument ball lies in $\{\re w > 0\}$ is such an inclusion function, with $U$ the set on which every divisor is nonzero and every logarithm and square-root argument has positive real part, by induction over the expression. (Arb returns a non-finite ball for a quotient whose divisor ball contains $0$; Arb's own complex logarithm does not signal a ball that crosses the branch cut, which is why the guards are needed.) Let $g = f \circ p$, $c_k = \frac{1}{2\pi}\int_0^{2\pi} g(\theta) e^{-ik\theta}\,d\theta$, and $\Sigma_\rho = \{\theta \in \C : |\im\theta| \le \rho\}$.

\begin{lemma}[strip cover]\label{lem:strip}
Cover the rectangle $R = [0, 2\pi] \times [-\rho, \rho]$ by finitely many closed rectangles $B$ with exact rational corners. For each $B$ let $\Theta_B$ be a complex ball containing $B$, $P_B = \sum_m a_m \exp(i m \Theta_B)$ evaluated in ball arithmetic, and $Y_B = \mathsf F(P_B)$. If every $Y_B$ is finite, then $p(\Sigma_\rho) \subset U$, $g$ is holomorphic on an open neighbourhood of $\Sigma_\rho$, and $\sup_{\Sigma_\rho} |g_i| \le S_i := \max_B \overline{|Y_{B,i}|}$, where $\overline{|\cdot|}$ is an upper bound of the modulus.
\end{lemma}

\begin{proof}
Every $\theta \in R$ lies in some $B$, so $p(\theta) \in P_B \subset U$ by the inclusion property. Since $p$ is $2\pi$-periodic, $p(\Sigma_\rho) = p(R) \subset U$; as $U$ is open and $p$ continuous, $p^{-1}(U)$ is an open set containing $\Sigma_\rho$, on which $g$ is holomorphic. Finally $|g_i(\theta)| \le \overline{|Y_{B,i}|} \le S_i$.
\end{proof}

The whole rectangle is covered, not only its edges: a pole of $g$ strictly inside the strip can coexist with a finite supremum on the edges, and then the next lemma fails. The programs refine the cover adaptively and optionally use a centred (mean-value) form on each box; both only affect how tight $S$ is.

\begin{lemma}[Cauchy estimate]\label{lem:cauchy}
If $g$ is holomorphic on an open set containing $\Sigma_\rho$, $2\pi$-periodic, and $|g| \le S$ on $\Sigma_\rho$, then $|c_m| \le S e^{-\rho|m|}$ for every $m \in \Z$.
\end{lemma}

\begin{proof}
For $m > 0$ integrate $g(\theta)e^{-im\theta}$ over the boundary of the rectangle with corners $0, 2\pi, 2\pi - i\rho, -i\rho$. The integral vanishes by Cauchy's theorem and the vertical sides cancel by periodicity, so $c_m = \frac{1}{2\pi}\int_0^{2\pi} g(x - i\rho)e^{-im(x - i\rho)}\,dx$ and $|e^{-im(x - i\rho)}| = e^{-m\rho}$. For $m < 0$ shift to $\im\theta = \rho$; $|c_0| \le S$ directly.
\end{proof}

\begin{lemma}[aliasing]\label{lem:alias}
Under the hypotheses of Lemma~\ref{lem:cauchy}, with $\hat c_k = \frac1M \sum_{j=0}^{M-1} g(2\pi j/M) e^{-2\pi ijk/M}$ and $|k| < M$,
\[ |\hat c_k - c_k| \le E_k := S\,\frac{e^{-\rho(M - k)} + e^{-\rho(M + k)}}{1 - e^{-\rho M}}. \]
\end{lemma}

\begin{proof}
By Lemma~\ref{lem:cauchy} $\sum_m |c_m| < \infty$, so the Fourier series converges absolutely and uniformly to $g$. Substituting and using $\frac1M\sum_j e^{2\pi ij(m-k)/M} = 1$ if $m \equiv k \pmod M$ and $0$ otherwise gives $\hat c_k = \sum_{l\in\Z} c_{k + lM}$. For $l \ge 1$, $k + lM > 0$ and $\sum_{l \ge 1}|c_{k+lM}| \le S e^{-\rho(M+k)}/(1 - e^{-\rho M})$; for $l = -j \le -1$, $|k - jM| = jM - k$ and the sum is at most $S e^{-\rho(M-k)}/(1 - e^{-\rho M})$.
\end{proof}

\begin{proposition}[enclosure of the coefficients]\label{prop:enclosure}
Let the strip cover succeed with bounds $S_i$. Evaluate the nodes $e^{im\theta_j}$, $\theta_j = 2\pi j/M$, as balls ($e^{i\pi q}$ for exact rationals $q$), $P_j = \sum_m a_m e^{im\theta_j}$ and $G_j = \mathsf F(P_j)$, and let $C_k = \frac1M\sum_j G_j e^{-ik\theta_j} + E_k\,([-1,1] + i[-1,1])$. Then for every $p$ with coefficients in the balls, every component $i$ and every $|k| \le K' < M$, $c_{i,k}[f \circ p] \in C_{i,k}$, and $|c_{i,m}| \le S_i e^{-\rho|m|}$ for every $m$.
\end{proposition}

\begin{proof}
The ball sum contains $\hat c_k$, and a complex number $w$ with $|w| \le E_k$ lies in the square $E_k([-1,1] + i[-1,1])$. Lemmas~\ref{lem:strip} to~\ref{lem:alias}.
\end{proof}

For $\nu > 0$ with $q = \nu e^{-\rho} < 1$ these bounds give the weighted tail $\sum_{|m| > K'} |c_m|\nu^{|m|} \le 2S q^{K'+1}/(1 - q)$.

\subsection{The zero-finding problem}\label{sec:F}

Fix $K = 32$ and a centre $\bar x = (\bar\omega, \bar a)$ of exact dyadic numbers, with $\bar a_m = 0$ for $|m| > K$, $\bar a_{-m} = \overline{\bar a_m}$ and $\bar\omega$ real; it is computed by an untrusted Newton iteration (\file{code/fourier/centre.py}) and stored in \file{code/fourier/data/centre_N*_K32.json}. Let $\nu = e^{1/4}$, let $\ell^1_\nu$ be the sequences $(b_m)_{m\in\Z}$ with $\|b\|_\nu = \sum_m |b_m|\nu^{|m|} < \infty$, a Banach algebra under convolution, and let $X = \C \times (\ell^1_\nu)^{18}$ with $\|x\| = \max(|\omega|, \max_k \|a_k\|_\nu)$ (the programs allow weights; all runs used weight $1$). Writing $\phi_a(\theta) = \sum_m a_m e^{im\theta}$, $g(a) = f \circ \phi_a$ and $g_m(a)$ for its Fourier coefficients, define $F$, on the open set of $x \in X$ for which $\phi_a$ maps the closed strip $|\im\theta| \le 1/4$ into the domain $U$ (it contains the ball $B_{r_*}(\bar x)$ used below, Section~\ref{sec:Z2}), with values in $X' := \C \times (\ell'_\nu)^{18}$, by
\begin{equation}\label{eq:F}
  F_{\mathrm{ph}}(x) = \sum_m a_{m,V} - s/\sigma_V, \qquad F_m(x) = i\omega m a_m - g_m(a) + d_m E a_m \quad (m \in \Z),
\end{equation}
where $\sigma_V = 2^{-2}$ is the scale of $V$, so $F_{\mathrm{ph}} = 0$ says $\phi_{a,V}(0) = s$ in physical units. Here $\ell'_\nu$ is the space of sequences with $\|b\|' = \sum_m |b_m|\nu^{|m|}/(1 + |m|) < \infty$, and $X'$ has the norm $\max(|\cdot|, \max_k\|\cdot\|')$; $F$ maps $X$ into $X'$ because $(i\omega ma_m)_m \in \ell'_\nu$ when $a \in \ell^1_\nu$, and $g(a)$, $d_mEa_m \in \ell^1_\nu \subset \ell'_\nu$. A zero of $F$ with $\omega > 0$ and $\phi_a$ real and $C^1$ solves \eqref{eq:profile}, and the phase condition removes the time-shift invariance.

\subsection{The radii-polynomial theorem}\label{sec:rp}

Let $A: X' \to X$ be an injective bounded linear operator, defined below, and $\Psi(x) = x - AF(x)$, a map from $X$ to $X$. Its fixed points are exactly the zeros of $F$.

\begin{lemma}[radii polynomial]\label{lem:rp}
Let $r_* > 0$ be such that $\Psi$ is defined and $C^1$ on the closed ball $B_{r_*}(\bar x)$, and suppose that
\[ \|AF(\bar x)\| \le Y_0, \qquad \|I - A\,DF(\bar x)\| \le Z_1, \qquad \|A(DF(x) - DF(\bar x))\| \le Z_2 \|x - \bar x\| \ \ (x \in B_{r_*}(\bar x)). \]
If $0 < r \le r_*$, $p(r) := Y_0 + (Z_1 - 1)r + Z_2 r^2/2 < 0$ and $Z_1 + Z_2 r < 1$, then $\Psi$ maps $B_r(\bar x)$ into itself and is a contraction there, so $F$ has exactly one zero in $B_r(\bar x)$.
\end{lemma}

\begin{proof}
$\|D\Psi(x)\| \le Z_1 + Z_2\|x - \bar x\|$. For $x \in B_r(\bar x)$, $\|\Psi(x) - \bar x\| \le \|\Psi(\bar x) - \bar x\| + \|\int_0^1 D\Psi(\bar x + t(x - \bar x))(x - \bar x)\,dt\| \le Y_0 + Z_1 r + Z_2 r^2/2 < r$. The mean value inequality on the convex ball gives the contraction constant $Z_1 + Z_2 r < 1$, and Banach's fixed point theorem gives one fixed point, which is a zero of $F$ since $A$ is injective.
\end{proof}

The programs check $p(r) < 0$ and $Z_1 + Z_2 r < 1$ in Arb at two exact dyadic radii: the small radius $r_{\mathrm{ex}}$ of Table~\ref{tab:existence}, which places the zero, and $r_{\mathrm{un}}$, the largest double not above $r_* = 10^{-12}$, which gives uniqueness (the contraction constant there is about $0.267$). The validity radius $r_*$ is chosen by hand, so $r_{\mathrm{un}}$ is not an optimal uniqueness radius.

\subsection{The operator \texorpdfstring{$A$}{A} and its tail}\label{sec:A}

Let $J(\theta) = Df(\phi_{\bar a}(\theta)) = \sum_n J_n e^{in\theta}$. The derivative $DF(\bar x)y$, $y = (y_\omega, (y_m))$, has phase component $\sum_m y_{m,V}$ and row $m$ equal to $iy_\omega m\bar a_m + i\bar\omega m y_m - \sum_n J_n y_{m-n} + d_m E y_m$. The operator $A$ is block diagonal with respect to the finite part (the phase row and modes $|m| \le K$, with the unknown $\omega$) and each tail mode $|m| > K$:
\begin{itemize}
\item $A_{\mathrm{fin}}$, an exact complex matrix of size $1 + 18(2K + 1)$, the floating-point inverse of the midpoint of the Galerkin matrix $J_{\mathrm{fin}}$ (the restriction of $DF(\bar x)$ to these rows and columns);
\item $A_m = (i\bar\omega m I - \hat J_0 + d_m E)^{-1}$ for $|m| > K$, where $\hat J_0$ is the exact real midpoint of the enclosure of $J_0$.
\end{itemize}
$A$ acts on $X'$ block by block. It is bounded from $X'$ to $X$ because $\sup_{|m| > K}(1 + |m|)|A_m| < \infty$ (Lemma~\ref{lem:Atail}), and it is injective because every block is: every $A_m$ is an inverse, and $A_{\mathrm{fin}}$ is invertible, which follows from $\|I - A_{\mathrm{fin}}J_{\mathrm{fin}}\| \le Z_1 < 1$, since the compression of $I - A\,DF(\bar x)$ to the finite modes is $I - A_{\mathrm{fin}}J_{\mathrm{fin}}$ and the coordinate projection onto the finite modes has norm 1 in $X$. Since $\hat J_0$ and $d_m$ are real, $A_{-m} = \overline{A_m}$, and only $m > 0$ is computed.

\begin{lemma}[tail of $A$]\label{lem:Atail}
For $K < m \le m_{\max}$ the matrices $A_m$ are enclosed by Arb inversion of a ball matrix whose only non-exact entries are the narrow balls of $d_m$ (exact for $N = 1$); Arb raises unless every matrix in the ball is certainly invertible, and otherwise returns a ball containing every inverse. Let $G_0 = |\hat J_0| + 4c\,E$ entrywise, $\Upsilon = \bar\omega(m_{\max} + 1)$, $G = G_0/\Upsilon$ and $\vartheta = \|G\|_\infty < 1$. Then for every $m > m_{\max}$, entrywise,
\[ |A_m| \le \frac{1}{\Upsilon} S_G, \qquad |mA_m| \le \frac{1}{\bar\omega} S_G, \qquad S_G := I + G + G^2 + \frac{\vartheta^3}{1 - \vartheta}\,\one\one^{\mathsf T}. \]
\end{lemma}

\begin{proof}
With $y = \bar\omega m \ge \Upsilon$ and $B_0 = \hat J_0 - d_m E$, $|B_0| \le G_0$ entrywise because $0 \le d_m \le 4c$. Then $A_m = (iy)^{-1}(I - B_0/(iy))^{-1} = (iy)^{-1}\sum_k (B_0/(iy))^k$, and entrywise $|A_m| \le y^{-1}\sum_k (G_0/y)^k \le \Upsilon^{-1}\sum_k G^k$. Each entry of $G^k$ is at most the corresponding row sum of $G^k$, which is at most $\vartheta^k$; summing over $k \ge 3$ gives the last term. Finally $|mA_m| = (y/\bar\omega)|A_m| \le \bar\omega^{-1}\sum_k (G_0/y)^k$.
\end{proof}

The program increases $m_{\max}$ until $\vartheta \le 1/2$ is certified ($m_{\max} = 391$ at $N = 8$). The suprema $\bar A_0 = \sup_{|m| > K}|A_m|$, $\bar A_1 = \sup_{|m| > K}|mA_m|$ and $C_n = \sup_{|m| > K}|A_m J'_n|$, with $J'_n = J_n$ for $n \ne 0$ and $J'_0 = J_0 - \hat J_0$, are taken entrywise over both ranges; for $|n| \le 12$, $C_n$ is the entrywise maximum of the explicit products $|A_mJ'_n|$ over $K < m \le m_{\max}$ and of the bound $\Upsilon^{-1}S_G|J'_n| \ge |A_m||J'_n|$ valid for every $m > m_{\max}$; for $12 < |n| \le K'$, $C_n = \bar A_0|J'_n|$. Negative $m$ are handled by $A_{-m}J'_n = \overline{A_m \overline{J'_n}}$.

\subsection{The bounds \texorpdfstring{$Y_0$}{Y0} and \texorpdfstring{$Z_1$}{Z1}}\label{sec:Y0Z1}

The Fourier data come from Proposition~\ref{prop:enclosure}: enclosures of $g_m(\bar a)$ and $J_n$ for $|m|, |n| \le K' = 2K + 16 = 80$ from an $M = 192$ node transform, with strip suprema $S_g$ and (entrywise) $S_J$ on $|\im\theta| \le \rho = 3/2$, and beyond $K'$ the bounds $|g_{k,m}| \le S_{g,k}e^{-\rho|m|}$, $|J_{n,jk}| \le S_{J,jk}e^{-\rho|n|}$. Let $q = \nu e^{-\rho} < 1$. For a linear operator $B$ on $X$ with blocks $B_{cc'}$, $\|B\| \le \max_c \sum_{c'}\|B_{cc'}\|$ with $\|B_{cc'}\|_{\ell^1_\nu \to \ell^1_\nu} = \sup_{m'}\sum_m |B(m,m')|\nu^{|m|}/\nu^{|m'|}$ (the column supremum).

\paragraph{$Y_0$.} The finite part of $AF(\bar x)$ is $A_{\mathrm{fin}}F_{\mathrm{fin}}(\bar x)$, computed in Arb. For $|m| > K$, $\bar a_m = 0$ and $F_m(\bar x) = -g_m$, so $(AF)_m = -A_m g_m$. For $K < |m| \le K'$ the computed enclosures of $A_m$ are used (all of them are explicit, since $m_{\max} = 391 > K' = 80$), and for $|m| > K'$, entrywise $|A_mg_m| \le \bar A_0 S_g e^{-\rho|m|}$, so $\sum_{|m| > K'}|(A_mg_m)_c|\nu^{|m|} \le \sum_k(\bar A_0)_{ck}S_{g,k}\,2q^{K'+1}/(1 - q)$. Then $Y_0 = \max_c\big(\|(AF)_{c,\mathrm{fin}}\|_\nu + \|(AF)_{c,\mathrm{tail}}\|_\nu\big)$, the phase/$\omega$ component counting as a component with the single mode $0$.

\paragraph{$Z_1$.} Let $B = I - A\,DF(\bar x)$. Finite rows and finite columns give $I - A_{\mathrm{fin}}J_{\mathrm{fin}}$, an Arb matrix product. A tail column $(k, m')$, $|m'| > K$, meets the finite rows through $-A_{\mathrm{fin}}w$, where $w$ has phase entry $1$ if $k = V$ and $0$ otherwise (the phase functional sums over all modes) and row $(j, m)$ entry $-J_{m-m',jk}$ (the terms $i\bar\omega m$ and $d_m$ vanish off the diagonal $m = m'$); for $K < |m'| \le K + 16$ the enclosures of $J$ are used, and for $|m'| \ge K + 17$ the bound $|J_{m-m',jk}| \le S_{J,jk}e^{-\rho(|m'| - \operatorname{sgn}(m')m)}$ makes every term of the weighted column sum divided by $\nu^{|m'|}$ a constant times $\nu^{-|m'|}$ or $(e^{-\rho}/\nu)^{|m'|}$, both decreasing in $|m'|$, so the value at $|m'| = K + 17$ bounds the supremum. For the tail rows, since $A_m(i\bar\omega m - \hat J_0 + d_m E) = I$,
\[ (By)_m = A_m\Big[(J_0 - \hat J_0)y_m + \sum_{n\ne0} J_n y_{m-n}\Big] = \sum_n A_m J'_n y_{m-n} \qquad (|m| > K), \]
($y_\omega$ does not enter because $\bar a_m = 0$ there), and $\nu^{|m|} \le \nu^{|m-n|}\nu^{|n|}$ gives $\|(By)_{c,\mathrm{tail}}\|_\nu \le \sum_k \mathcal T_{ck}\|y_k\|_\nu$ with $\mathcal T = \sum_{|n| \le K'} C_n\nu^{|n|} + \bar A_0 S_J\, 2q^{K'+1}/(1 - q)$, for every column. For an output component $c$ and an input component $c'$ let $f\!f_{cc'}$ be the largest weighted column sum of the finite rows of block $(c, c')$ over the finite columns, and $f\!t_{cc'}$ the same over the tail columns (bounded as above). A column of $B$ meets the finite rows and the tail rows, so its weighted sum in block $(c,c')$ is at most $\max(f\!f_{cc'}, f\!t_{cc'}) + \mathcal T_{cc'}$, and
\[ Z_1 = \max_c\sum_{c'}\big(\max(f\!f_{cc'}, f\!t_{cc'}) + \mathcal T_{cc'}\big). \]
The $\omega$ column has a single finite column and no tail rows ($iy_\omega m\bar a_m = 0$ for $|m| > K$), and the $\omega$ row of $B$ (the output of the phase column of $A_{\mathrm{fin}}$) has no tail part, since $A$ has no tail in that row, so for it only the finite-row terms enter.

\subsection{The bound \texorpdfstring{$Z_2$}{Z2} by a polydisc majorant}\label{sec:Z2}

Let $R_k = 1/1024$ for every component, and let $M_k$ be the strip supremum (Lemma~\ref{lem:strip}, on $|\im\theta| \le \rho_2 = 1$) of $f_k$ over the family of trigonometric polynomials obtained from $\phi_{\bar a}$ by replacing each constant coefficient $\bar a_{0,j}$ by the complex box of half-width $R_j$ about it. Since that box contains the disc of radius $R_j$, $M_k \ge \sup\{|f_k(\phi_{\bar a}(\theta) + w)| : |\im\theta| \le 1, |w_j| \le R_j\}$, and the finite cover certifies that $f$ is holomorphic on an open set containing these points.

\begin{lemma}\label{lem:Z2}
Let $q_2 = \nu e^{-\rho_2} < 1$ and $Q_2 = (1 + q_2)/(1 - q_2)$. For $h \in (\ell^1_\nu)^{18}$ with $\|h_j\|_\nu < R_j$, the map $a \mapsto g_k(a)$ is analytic at $\bar a + h$ and
\[ \|D^2 g_k(\bar a + h)[u, v]\|_\nu \le \sum_{j,l}\partial_j\partial_l\Phi_k(\|h\|)\,\|u_j\|_\nu\|v_l\|_\nu, \qquad \Phi_k(t) = M_k Q_2\prod_j (1 - t_j/R_j)^{-1}. \]
\end{lemma}

\begin{proof}
For $|\im\theta| \le 1$ the Taylor coefficients $c_\alpha(\theta) = \partial^\alpha f_k(\phi_{\bar a}(\theta))/\alpha!$ satisfy $|c_\alpha(\theta)| \le M_k R^{-\alpha}$ by Cauchy's inequality on the closed polydisc, and they are holomorphic and $2\pi$-periodic in $\theta$; Lemma~\ref{lem:cauchy} gives $|[c_\alpha]_m| \le M_k R^{-\alpha}e^{-|m|}$ and so $\|c_\alpha\|_\nu \le M_k R^{-\alpha}Q_2$. In the Banach algebra $\ell^1_\nu$ the series $\sum_\alpha c_\alpha h^\alpha$ (convolution powers) converges absolutely, dominated by $\Phi_k(\|h\|)$. For real $\theta$, $|h_j(\theta)| \le \|h_j\|_\nu < R_j$, so the Taylor series of $f_k$ at $\phi_{\bar a}(\theta)$ converges to $f_k(\phi_{\bar a}(\theta) + h(\theta))$; convergence in $\ell^1_\nu$ is uniform convergence, so the two functions agree and $g_k(\bar a + h) = \sum_\alpha c_\alpha h^\alpha$. Differentiating termwise and comparing with the majorant, whose coefficients are nonnegative, gives the bound.
\end{proof}

For $x = \bar x + \Delta$ with $\|\Delta\| \le r$ and $\|y\| \le 1$,
\[ (DF(x) - DF(\bar x))y = im(y_\omega\,\Delta a_m + \Delta\omega\,y_m) - (Dg(\bar a + \Delta a) - Dg(\bar a))y_a. \]
The first term $b = y_\omega\Delta a + \Delta\omega\,y_a$ has $\|b_k\|_\nu \le 2r$, and $\|A(imb)\|_c \le \sum_k (N_1 + \bar A_1)_{ck}\|b_k\|$, with $N_1$ the block norms of $A_{\mathrm{fin}}\operatorname{diag}(im)$. The second equals $\int_0^1 D^2g(\bar a + t\Delta a)[\Delta a, y_a]\,dt$, whose component $k$ has norm at most $rW_k(r)$ with
\[ W_k(r) = M_kQ_2P(r)\Big[\Big(\sum_j s_j\Big)^2 + \sum_j s_j^2\Big], \qquad s_j = \frac{1}{R_j - r}, \quad P(r) = \prod_j\frac{R_j}{R_j - r}, \]
because $\partial_jP = P/(R_j - t_j)$, $\partial_l\partial_jP = P/((R_j - t_j)(R_l - t_l))$ for $j \ne l$ and $2P/(R_j - t_j)^2$ for $j = l$, and $\Phi_k$ is increasing. Hence
\[ Z_2 = \max_c\sum_k\big[2(N_1 + \bar A_1)_{ck} + (N_0 + \bar A_0)_{ck}W_k(r_*)\big], \]
with $N_0$ the block norms of $A_{\mathrm{fin}}$ on the coefficient columns; in the $\omega$ row only $N_0$ and $N_1$ enter. $\Psi$ is defined and $C^1$ on $B_{r_*}(\bar x)$ because $r_* < R_j$, $g$ is analytic there, the term $A(i\omega ma_m)$ is a bounded bilinear map ($\sup_{|m| > K}|m||A_m| < \infty$ by Lemma~\ref{lem:Atail}), and the other terms are bounded linear.

\subsection{The real solution, its period and its wave number}\label{sec:real}

\begin{lemma}\label{lem:real}
Let the hypotheses of Lemma~\ref{lem:rp} hold at $r_{\mathrm{ex}}$ and $r_{\mathrm{un}}$, and let $x_* = (\omega_*, a_*)$ be the zero. Then $\omega_*$ is real and positive, $\phi_* = \phi_{a_*}$ is real-valued, analytic on $|\im\theta| < 1/4$, and solves \eqref{eq:profile}. If also $|\bar a_{1,V}| > r_{\mathrm{ex}}/\nu$, then $\phi_*$ has minimal period $2\pi$, and for $N \ge 2$ the ring solution $x_j(t) = \phi_*(\omega_*t + 2\pi j/N)$ has minimal period $T = 2\pi/\omega_*$, is not synchronous and is a rotating 1-wave in the sense of Theorem~\ref{thm:ring}(b).
\end{lemma}

\begin{proof}
Let $\kappa(\omega, a) = (\operatorname{conj}\omega, (\operatorname{conj} a_{-m})_m)$, an isometry of $X$ with $\kappa(\bar x) = \bar x$. The program \file{code/fourier/tp06_18d_arb.py} that evaluates $f$ is a fixed composition of $+, -, \times, /$, integer powers, 55 exponentials, 4 logarithms and 3 square roots with real constants, and contains no absolute value, real or imaginary part, comparison, branch or conversion. Each of these operations commutes with complex conjugation wherever it is defined: the arithmetic operations algebraically, $\exp$ because its power series has real coefficients, and the principal logarithm and square root on $\{\re w > 0\}$, the only arguments admitted, because that set is invariant under conjugation and $\operatorname{Log}\bar w = \overline{\operatorname{Log} w}$, $\sqrt{\bar w} = \overline{\sqrt w}$ there. By induction over the expression, if $z$ lies in the domain then so does $\bar z$, and $f(\bar z) = \overline{f(z)}$. The points used are $\phi_{\bar a}(\theta) + h(\theta)$ with $\theta$ real and $|h_j(\theta)| < R_j$; this family is closed under conjugation ($\phi_{\bar a}(\theta)$ is real) and lies in the certified domain. Hence $g_m(\kappa a) = \overline{g_{-m}(a)}$, and since $d_{-m} = d_m$ and $s/\sigma_V$ are real, $F(\kappa x) = \overline{F(x)}$ in the corresponding sense; so $\kappa$ maps zeros of $F$ in $B_{r_{\mathrm{un}}}(\bar x)$ to zeros in the same ball, and uniqueness gives $\kappa x_* = x_*$: $\omega_*$ is real and $\phi_*$ is real-valued. As $a_* \in (\ell^1_\nu)^{18}$, $\phi_*$ is analytic on $|\im\theta| < 1/4$; since $i\omega_* m a_{*,m} = g_m - d_m E a_{*,m}$, the coefficient sequence of $\phi_*'$ is in $\ell^1_\nu$, and \eqref{eq:profile} holds pointwise. $\omega_* \ge \bar\omega - r_{\mathrm{ex}} > 0$ is checked in Arb.

Since $\|a_{*,V} - \bar a_V\|_\nu \ge \nu|a_{*,1,V} - \bar a_{1,V}|$, $|a_{*,1,V} - \bar a_{1,V}| \le r_{\mathrm{ex}}/\nu$, so $a_{*,1,V} \ne 0$. A function of minimal period $2\pi/k$, $k \ge 2$, has $a_m = 0$ unless $k \mid m$; hence $\phi_*$ has minimal period $2\pi$ and each cell has minimal period $T$. A synchronous solution would make $\phi_*$ periodic with period $2\pi/N$. If $x_j(t) = x_0(t + jkT/N)$ for all $j$ and some $k \not\equiv 1 \pmod N$, then with $j = 1$, $\phi_*(\theta + 2\pi/N) = \phi_*(\theta + 2\pi k/N)$, so $\phi_*$ would have the period $2\pi(k-1)/N \not\equiv 0 \pmod{2\pi}$, contradicting minimality.
\end{proof}

\begin{proof}[Proof of Theorem~\ref{thm:ring}(a), (b) and the existence part of Theorem~\ref{thm:cell}(ii)]
For each $N \in \{1, 8, 16, 32, 64\}$ the program \file{code/fourier/existence.py} computes, in Arb, the strip covers (all leaves finite, full strip), the enclosures of Proposition~\ref{prop:enclosure}, the tail of Lemma~\ref{lem:Atail}, $Y_0$, $Z_1$ and $Z_2$ as above, checks $p(r) < 0$ and $Z_1 + Z_2r < 1$ at $r = r_{\mathrm{ex}}$ and at $r = r_{\mathrm{un}} \le r_*$, checks $\bar\omega - r_{\mathrm{ex}} > 0$, and $|\bar a_{1,V}| > r_{\mathrm{ex}}/\nu$, and writes the interval $2\pi/[\omega_*]$ with outward-rounded decimal ends. The records \file{data/fourier-existence-N*.json} store every bound; for example at $N = 8$: $Y_0 = 1.3096 \times 10^{-28}$, $Z_1 = 0.20339$, $Z_2 \le 6.4002 \times 10^{10}$, $|\bar a_{1,V}| \ge 0.11524$. Lemmas~\ref{lem:rp} and~\ref{lem:real} give the statements.
\end{proof}

\section{Stability through one Hill operator}\label{sec:stability}

\subsection{Setting}\label{sec:hill}

Let $\phi$, $\omega$, $T$, $\tau$ be as in Section~\ref{sec:ringsym}, with $\phi$ solving \eqref{eq:profile}, and let $A(\theta) = Df(\phi(\theta)) = \sum_n A_n e^{in\theta}$, a real analytic $2\pi$-periodic matrix function, so $A_{-n} = \overline{A_n}$ and $\|A_n\|$ decays exponentially. The linearization along $x(t)$ is $y' = L(t)y$ with $(L(t)y)_j = A(\theta_j(t))y_j + cE(y_{j-1} - 2y_j + y_{j+1})$; let $Y(t)$ be its fundamental matrix with $Y(0) = I$, and $M_\tau = Q^{-1}Y(\tau)$, the derivative at $x(0)$ of $h = Q^{-1}\varphi_\tau$. The Floquet multipliers are the eigenvalues of $Y(T)$, with algebraic multiplicity. For a matrix $M$ let $m(\lambda; M)$ be the algebraic multiplicity and $G_\lambda(M)$ the generalized eigenspace; for an operator $H$, $G_\mu(H) = \bigcup_k\ker(H - \mu)^k$ and $m(\mu; H) = \dim G_\mu(H)$.

Let $|\cdot|_1$ be the $\ell^1$ norm on $\C^{18}$, $\|B\|_{1\to1}$ the induced norm, $\Xsp = \ell^1(\Z; \C^{18})$ with $\|P\| = \sum_m|P_m|_1$, and $\Dsp = \{P \in \Xsp : \sum_m |m||P_m|_1 < \infty\}$. A sequence $P \in \Dsp$ is the coefficient sequence of the $C^1$ function $p(\theta) = \sum_m P_m e^{im\theta}$. The Hill operator $H_0: \Dsp \to \Xsp$ is
\begin{equation}\label{eq:H0}
  (H_0P)_m = -i\omega mP_m + \sum_n A_nP_{m-n} - d_mEP_m,
\end{equation}
the coefficient sequence of $\Lop p := -\omega p' + Ap + cE\Delta_Np$. The ansatz $y_j(t) = e^{\mu t}p(\theta_j(t))$ solves the linearization if and only if $\Lop p = \mu p$. (A numerical toy model with three-dimensional cells and $N = 5$ reproduced $e^{\mu\tau} \in \spec M_\tau$ to $3 \times 10^{-13}$ with this sign convention and missed by $4 \times 10^{-2}$ with the opposite sign of the diagonal; this check is not part of any proof.)

We use three standard facts about a closed operator $H$ with compact resolvent, from Kato's book \cite[Ch.~III, Sect.~6]{kato1976}: its spectrum consists of isolated eigenvalues of finite algebraic multiplicity; for a bounded open set $\Omega$ whose boundary $\Gamma$ lies in the resolvent set, the Riesz projection $\frac{1}{2\pi i}\oint_\Gamma(\mu - H)^{-1}d\mu$ ($\Gamma$ counterclockwise) is a projection onto the sum of the generalized eigenspaces of the eigenvalues in $\Omega$, so its rank is $n(H, \Omega) := \sum_{\mu\in\Omega}m(\mu; H)$; and the resolvent is holomorphic on the resolvent set. (We cite these by chapter and section only; their statements and numbering in Kato's book have not been checked against a copy.)

\begin{lemma}\label{lem:H0}
$H_0$ is closed on $\Dsp$ and has compact resolvent, and $\{\re\mu > \alpha\}$ lies in its resolvent set, where $\alpha = \sum_n\|A_n\|_{1\to1} + 4c$.
\end{lemma}

\begin{proof}
Write $H_0 = D_0 + B$ with $(D_0P)_m = -i\omega mP_m$ on $\Dsp$ and $B$ the rest. By Young's inequality and $d_m \le 4c$, $\|B\| \le \alpha$. $D_0$ is closed (a diagonal operator on its maximal domain), and for $\lambda \notin i\omega\Z$, $(\lambda - D_0)^{-1}$ is diagonal with entries $(\lambda + i\omega m)^{-1}$, compact as the norm limit of its finite truncations. If $\re\lambda > \alpha$, then $\|(\lambda - D_0)^{-1}\| \le 1/\re\lambda$, so $\|B(\lambda - D_0)^{-1}\| < 1$ and $\lambda - H_0 = (I - B(\lambda - D_0)^{-1})(\lambda - D_0)$ is a bijection $\Dsp \to \Xsp$ with compact inverse. $H_0$ is closed as a closed operator plus a bounded one, and compactness of the resolvent at one point gives it everywhere by the resolvent identity $R(\lambda) = R(\lambda_0)(I + (\lambda_0 - \lambda)R(\lambda))$.
\end{proof}

\subsection{The Hill-sector theorem}\label{sec:hillsector}

\begin{lemma}[twisted periodicity]\label{lem:twist}
$L(t + \tau) = QL(t)Q^{-1}$, $Y(t + \tau) = QY(t)M_\tau$, $Y(k\tau) = Q^kM_\tau^k$ and $Y(T) = M_\tau^N$.
\end{lemma}

\begin{proof}
$(QL(t)Q^{-1}y)_j = A(\theta_{j+1}(t))y_j + cE(y_{j-1} - 2y_j + y_{j+1}) = (L(t + \tau)y)_j$, because $(Q^{-1}y)_k = y_{k-1}$ and $\theta_{j+1}(t) = \theta_j(t + \tau)$. So $Z(t) = Q^{-1}Y(t + \tau)$ solves $Z' = L(t)Z$ with $Z(0) = M_\tau$, and $Z(t) = Y(t)M_\tau$. Induction gives $Y(k\tau) = Q^kM_\tau^k$, and $Q^N = I$ gives $Y(T) = M_\tau^N$.
\end{proof}

\begin{thmx}[Hill sectors]\label{thm:hill}
For every $\mu \in \C$, $m(\mu; H_0) = m(e^{\mu\tau}; M_\tau)$.
\end{thmx}

\begin{proof}
Let $\lambda = e^{\mu\tau}$.

\emph{$m(\mu; H_0) \le m(\lambda; M_\tau)$.} Let $G = G_\mu(H_0)$, read through the Fourier correspondence as a finite-dimensional (Lemma~\ref{lem:H0}) $\Lop$-invariant space of $C^1$ functions, with $\Lop|_G = \mu + K_G$, $K_G$ nilpotent. For $p \in G$ put $u(t) = e^{t\Lop_G}p \in G$ and $y_j(t) = u(t)(\theta_j(t))$. The chain rule gives $y_j' = (\Lop_Gu + \omega u')(\theta_j) = (Au + cE\Delta_Nu)(\theta_j) = A(\theta_j)y_j + cE(y_{j-1} - 2y_j + y_{j+1})$, so $y$ solves the linearization with $y(0) = \iota^*p$, $\iota^*p := (p(2\pi j/N))_{j\in\Z_N}$. As $\theta_j(t + \tau) = \theta_{j+1}(t)$, $y(\tau) = Q\iota^*(e^{\tau\Lop_G}p)$; also $y(\tau) = Y(\tau)\iota^*p$. Hence $M_\tau\iota^* = \iota^*e^{\tau\Lop_G}$ on $G$. The map $\iota^*$ is injective on $G$: if $\iota^*p = 0$ then $y = 0$, so for all $t$, $j$ and integers $n$, $0 = u(t + nT)(\theta_j(t)) = e^{\mu(t+nT)}\sum_k\frac{(t+nT)^k}{k!}(K_G^kp)(\theta_j(t))$. For fixed $t, j$ this is $e^{\mu(t+nT)}$ times a polynomial in $n$ vanishing at every integer, so all its coefficients vanish; if $p \ne 0$ and $k^*$ is the largest $k$ with $K_G^kp \ne 0$, the coefficient of $n^{k^*}$ is $\frac{T^{k^*}}{k^*!}(K_G^{k^*}p)(\theta_j(t))$, so $K_G^{k^*}p$ vanishes at every point, a contradiction. So $\iota^*(G)$ is an $M_\tau$-invariant space of dimension $\dim G$ on which $M_\tau$ is similar to $\lambda e^{\tau K_G}$, whose only eigenvalue is $\lambda$; hence $\iota^*(G) \subset G_\lambda(M_\tau)$.

\emph{$m(\mu; H_0) \ge m(\lambda; M_\tau)$.} Let $W = G_\lambda(M_\tau)$, $r = \dim W \ge 1$, and on $W$ write $M_\tau = \lambda(I + K)$, $K$ nilpotent. Put $B = \mu I + \tau^{-1}\sum_{k=1}^r(-1)^{k+1}K^k/k$ on $W$. Then $e^{\tau B} = \lambda\exp(\log(I + K)) = M_\tau$ on $W$ (an identity of formal power series, valid for nilpotent $K$), $B$ commutes with $M_\tau|_W$, and the only eigenvalue of $B$ is $\mu$. Let $\Pi(t) = Y(t)e^{-tB}: W \to \C^{18N}$. By Lemma~\ref{lem:twist}, $\Pi(t + \tau)w = QY(t)M_\tau e^{-\tau B}e^{-tB}w = Q\Pi(t)w$, so $\Pi_j(t)w = \Pi_0(t + j\tau)w$ and $\Pi_0(t + T)w = \Pi_0(t)w$. Hence $p_w(\theta) := \Pi_0(\theta/\omega)w$ is $2\pi$-periodic, $C^\infty$, with coefficients in $\Dsp$, and $\Pi_j(t)w = p_w(\theta_j(t))$. Differentiating $Y(t)w = \Pi(t)e^{tB}w$ gives $L(t)\Pi(t)w = \Pi'(t)w + \Pi(t)Bw$, whose component $0$ at $t = \theta/\omega$ reads $\Lop p_w = p_{Bw}$. The map $w \mapsto p_w$ is injective ($p_w(2\pi j/N) = w_j$), so its image is an $\Lop$-invariant space of dimension $r$ on which $(\Lop - \mu)^r = 0$; it lies in $G_\mu(H_0)$.
\end{proof}

\begin{corollary}\label{cor:strip}
(i) $\spec H_0 = \{\mu : e^{\mu\tau} \in \spec M_\tau\}$, invariant with multiplicities under $\mu \mapsto \mu + i\omega N$ and $\mu \mapsto \bar\mu$.
(ii) For every real $a$, on the half-open strip $S_a = \{a \le \im\mu < a + \omega N\}$ the map $\mu \mapsto e^{\mu\tau}$ is a bijection from $\spec H_0 \cap S_a$ onto $\spec M_\tau$ preserving algebraic multiplicities, and these multiplicities add up to $18N$.
(iii) For $\rho \ne 0$, $m(\rho; Y(T)) = \sum\{m(\mu; H_0) : \mu \in \spec H_0 \cap S_a,\ e^{\mu T} = \rho\}$.
(iv) Let $H_q$ be $H_0$ with $d_m$ replaced by $d_{m+q}$. Then $H_q$ is similar to $H_0 + i\omega q$ by the shift $(SP)_k = P_{k-q}$, so the ``$N$ sectors'' $H_q$, $q \in \Z_N$, carry no information beyond $H_0$.
\end{corollary}

\begin{proof}
(i) Theorem~\ref{thm:hill} in both directions (every logarithm of an eigenvalue of $M_\tau$ is an eigenvalue of $H_0$), $e^{i\omega N\tau} = 1$, and $M_\tau$ real. (ii) $e^{\mu\tau}$ has period $i\omega N$, so it is injective on $S_a$ and onto $\C\setminus\{0\}$; $M_\tau$ is invertible. (iii) $Y(T) = M_\tau^N$; $\C^{18N}$ is the sum of the $G_\lambda(M_\tau)$, each invariant under $M_\tau^N$ with only eigenvalue $\lambda^N = e^{\mu T}$. (iv) Substitute $k = m + q$.
\end{proof}

\begin{remark}[half-open strips]\label{rem:halfopen}
The count $18N$ holds in a half-open strip of height $\omega N$. A closed strip $|\im\mu| \le \omega N/2$ can count an eigenvalue twice: $M_\tau$ is real, so a negative real eigenvalue of $M_\tau$ persists under perturbation and its logarithms lie exactly on $\im\mu \in \omega N/2 + \omega N\Z$ (the toy model above shows such eigenvalues). An early draft of the argument used the closed strip; it was corrected to the half-open strip before any certificate was run, and the certificate below never relies on a count in a strip: it counts eigenvalues only inside a rectangle whose boundary is in the resolvent set.
\end{remark}

\begin{corollary}[section maps]\label{cor:section}
Let $\Sigma$ be a $C^1$ hypersurface through $x^* = x(0)$ transverse to $F(x^*)$, and $g(z) = Q^{-1}\varphi_{t_g(z)}(z)$ a section map with $t_g$ of class $C^1$, $t_g(x^*) = \tau$, $g(\Sigma) \subset \Sigma$. Then $\det(\lambda I - M_\tau) = (\lambda - 1)\det(\lambda I - Dg(x^*))$.
\end{corollary}

\begin{proof}
$Y(\tau)F(x^*) = F(x(\tau)) = QF(x^*)$, so $M_\tau F(x^*) = F(x^*)$. $Dg(x^*) = M_\tau + F(x^*)Dt_g(x^*)$ on $T\Sigma$ maps $T\Sigma$ into itself, so it is the compression of $M_\tau$ along $F(x^*)$; in the splitting $\operatorname{span}F(x^*) \oplus T\Sigma$, $M_\tau$ is block upper triangular with diagonal blocks $1$ and $Dg(x^*)$.
\end{proof}

\subsection{The trivial multiplier}\label{sec:trivial}

\begin{lemma}\label{lem:trivial}
(a) $P^* = (ima_m)_m$, the coefficients of $\phi'$, is a nonzero element of $\Dsp$ with $H_0P^* = 0$. (b) $m(1; Y(T)) = \sum_{q=0}^{N-1}m(i\omega q; H_0)$. (c) The multiplier $1$ is algebraically simple if and only if $m(0; H_0) = 1$ and $i\omega q \notin \spec H_0$ for $q = 1, \dots, N - 1$.
\end{lemma}

\begin{proof}
(a) $\phi$ is analytic and nonconstant; differentiating \eqref{eq:profile} gives $\Lop\phi' = 0$. (b) Corollary~\ref{cor:strip}(iii) with $\rho = 1$, $a = -\omega/2$: $e^{\mu T} = 1$ iff $\mu \in i\omega\Z$, and $i\omega q \in S_{-\omega/2}$ iff $q = 0, \dots, N-1$. (c) follows from (a) and (b).
\end{proof}

An eigenvalue $i\omega q$, $q \not\equiv 0 \pmod N$, would be a neutral mode of the reduced map (a spatial phase twist) that merges with the multiplier $1$ of $Y(T)$. The certificate below excludes it together with everything else in $\{\re\mu \ge -\delta\}$.

\subsection{Exclusion and counting by a Riesz-projection homotopy}\label{sec:riesz}

Goal: for a given $\delta > 0$, prove
\begin{equation}\label{eq:G}
  \spec H_0 \cap \{\re\mu \ge -\delta\} = i\omega N\Z, \quad\text{each eigenvalue algebraically simple.}
\end{equation}

\begin{lemma}\label{lem:projrank}
If $P$, $P'$ are bounded projections on a Banach space, $\|P - P'\| < 1$ and $\rank P < \infty$, then $\rank P' = \rank P$.
\end{lemma}

\begin{proof}
If $x \in \Ran P'$ and $Px = 0$ then $\|x\| = \|(P' - P)x\| < \|x\|$ unless $x = 0$; so $P$ maps $\Ran P'$ injectively into $\Ran P$. Exchange the roles.
\end{proof}

\begin{lemma}[homotopy count]\label{lem:homotopy}
Let $\hat D$ be closed with compact resolvent, $\hat E$ bounded, $\Omega$ a bounded open rectangle whose boundary $\Gamma$ lies in the resolvent set of $\hat D$, and suppose $I - s\hat E(\mu - \hat D)^{-1}$ is invertible for all $(s, \mu) \in [0,1] \times \Gamma$. Then, with $H(s) = \hat D + s\hat E$, $\Gamma$ lies in the resolvent set of every $H(s)$ and $n(H(1), \Omega) = n(\hat D, \Omega)$.
\end{lemma}

\begin{proof}
$\mu - H(s) = (I - s\hat E(\mu - \hat D)^{-1})(\mu - \hat D)$ is a bijection with compact inverse $R_s(\mu)$, norm continuous on the compact set $[0,1] \times \Gamma$. The spectrum of $H(s)$ is discrete (Lemma~\ref{lem:H0}'s argument), so $\Omega$ contains finitely many eigenvalues, and the Riesz projection $P(s)$ has rank $n(H(s), \Omega)$. $P(s)$ is norm continuous in $s$, so by Lemma~\ref{lem:projrank} its rank is locally constant, hence constant.
\end{proof}

\paragraph{The comparison operator.} The program fixes:
\begin{itemize}
\item cell coordinates: an exact diagonal $S = \operatorname{diag}(2^{k_1}, \dots, 2^{k_{18}})$, applied in every mode (its exponents $k_l$ are chosen per $N$ and stored in the records; they are not the scale exponents $e_i$ of Section~\ref{sec:cell}). This is a similarity; it commutes with $E$, leaves $-i\omega m$ and $d_m$ unchanged, and changes no spectrum or multiplicity. From here on $A_n$ and every norm are taken in these coordinates;
\item a window $W = \{|m| \le K_e\}$ ($n_W = 18(2K_e + 1)$ coordinates) and the tail $\mathrm{Tl} = \{|m| > K_e\}$;
\item an invertible complex $n_W \times n_W$ matrix $V$ and a diagonal $\Lambda = \operatorname{diag}(\lambda_j)$ with exactly representable entries (floating-point eigenvectors and eigenvalues of the window block);
\item an exact real $18 \times 18$ matrix $A_0^c$ (the midpoint of the enclosure of $A_0$), and for each residue $r \in \{0, \dots, \lfloor N/2\rfloor\}$ an invertible $U_r$ and an exact diagonal $\Lambda_r$ approximately diagonalizing $X_r := A_0^c - d_rE$;
\item the rectangle $\Omega = (-\delta, R_0) \times (a, b)$ (real part, imaginary part), with boundary $\Gamma$.
\end{itemize}
Let $\mathrm{Scal}$ be $V$ on the window coordinates and the identity on the tail; it is bounded, boundedly invertible and maps $\Dsp$ onto $\Dsp$. Define
\[ \hat D = \Lambda \oplus \bigoplus_{m\in\mathrm{Tl}}B_m, \qquad B_m = -i\omega mI + X_{m \bmod N}, \qquad \hat E = \mathrm{Scal}^{-1}H_0\,\mathrm{Scal} - \hat D. \]
In blocks, $\hat E_{WW} = R_W := V^{-1}H_{WW}V - \Lambda$, $\hat E_{WT} = V^{-1}H_{WT}$, $\hat E_{TW} = H_{TW}V$ and $(\hat E_{TT}P)_m = \sum_{n \ne 0,\,m-n\in\mathrm{Tl}}A_nP_{m-n} + (A_0 - A_0^c)P_m$. $\hat E$ is bounded. The exact $\omega$ stays in $B_m$: an error $(\omega - \bar\omega)m$ is unbounded in $m$ and cannot be moved into $\hat E$. The norm on the $\mathrm{Scal}$-coordinates $v$ is $\|v\| = \sum_j|v_j| + \sum_{m\in\mathrm{Tl}}|v_m|_1$ (all weights equal to $1$ in the runs), and the norm of an operator is the supremum of its column sums.

\paragraph{Notation for the bounds.} All are upper bounds unless compared from below:
\begin{itemize}
\item $\operatorname{dist}_j$: the distance from $\lambda_j$ to $\Gamma$ (a lower bound);
\item $f\!m_j$: the $\ell^1$ norm of column $j$ of $R_W$; $r_j = \sum_{w\in W}t_w|V_{(w,\cdot),j}|_1$ with $t_w = \sum_{|w+n| > K_e}\|A_n\|_{1\to1}$;
\item $b_m = \max_k\sum_j|(V^{-1}H_{W,m})_{j,k}|$ for $m \in \mathrm{Tl}$, where $H_{W,m}$ has entries $(A_{w-m})_{l,k}$ in row $(w, l)$; $\hat b = \sup_{m\in\mathrm{Tl}}b_m$;
\item $\sigma_{\mathrm{off}} = \sum_{n\ne0}\|A_n\|_{1\to1}$ and $\theta_c = \sigma_{\mathrm{off}} + \|A_0 - A_0^c\|_{1\to1}$;
\item $\rho_T$: a bound of $\|(\mu - B_m)^{-1}\|_{1\to1}$ over $\mu \in \overline\Omega$ and $m \in \mathrm{Tl}$ (Lemma~\ref{lem:tail}), and $\theta_T = \theta_c\rho_T$.
\end{itemize}

\begin{lemma}[tail blocks]\label{lem:tail}
Let $h = \max(|a|, |b|)$, $\omega_{\mathrm{lo}} \le \omega$, $g_0 = \omega_{\mathrm{lo}}(K_e + 1) - h$, $\eta > 0$, and
\[ \gamma_r = \min_l\max\big(-\delta - \eta - \re\lambda_{r,l},\ g_0 - \eta - |\im\lambda_{r,l}|\big), \quad F_r = U_r^{-1}X_rU_r - \Lambda_r, \quad \kappa_r \ge \|U_r\|_{1\to1}\|U_r^{-1}\|_{1\to1}. \]
If $\gamma_r > \|F_r\|_{1\to1}$ for every $r$, then for every $\mu$ in the open neighbourhood $N_\eta = \{-\delta - \eta < \re\mu < R_0 + \eta,\ |\im\mu| < h + \eta\}$ of $\overline\Omega$ and every $m \in \mathrm{Tl}$, $\mu - B_m$ is invertible with $\|(\mu - B_m)^{-1}\|_{1\to1} \le \rho_T := \max_r\kappa_r/(\gamma_r - \|F_r\|_{1\to1})$. Moreover, for fixed $\mu$, $\|(\mu - B_m)^{-1}\|_{1\to1} \to 0$ as $|m| \to \infty$.
\end{lemma}

\begin{proof}
$\mu - B_m = z - X_r$ with $z = \mu + i\omega m$ and $r = m \bmod N$ (or $N - r$, which gives the same $X_r$). For $m \ge K_e + 1$, $\re z = \re\mu > -\delta - \eta$ and $\im z \ge \omega_{\mathrm{lo}}(K_e + 1) - h - \eta = g_0 - \eta$; for $m \le -(K_e + 1)$ the same holds for $\bar z$, and since $X_r$ is real, $(\bar z - X_r)^{-1} = \overline{(z - X_r)^{-1}}$ has the same norm. For such $z$, $|z - \lambda_{r,l}| \ge \max(-\delta - \eta - \re\lambda_{r,l}, g_0 - \eta - |\im\lambda_{r,l}|) \ge \gamma_r$, so $\|(z - \Lambda_r)^{-1}\| \le 1/\gamma_r$, and $z - X_r = U_r(z - \Lambda_r - F_r)U_r^{-1}$ with the Neumann series gives the bound. The decay follows from the Neumann bound $1/(|\mu + i\omega m| - \|X_r\|)$.
\end{proof}

\begin{lemma}[comparison operator]\label{lem:compare}
If the hypothesis of Lemma~\ref{lem:tail} holds and $\operatorname{dist}_j > 0$ for every $j$, then $\hat D$ has compact resolvent, $\Gamma$ lies in its resolvent set, $n(\hat D, \Omega) = \#\{j : \lambda_j \in \Omega\}$, and $\|(\mu - \hat D)^{-1}\| \le \max(\max_j 1/\operatorname{dist}_j, \rho_T)$ on $\Gamma$.
\end{lemma}

\begin{proof}
$\hat D$ is block diagonal. By Lemma~\ref{lem:tail}, for the tail part $D_T$ the operator $\mu - D_T$ has a bounded inverse for every $\mu \in N_\eta$, with norm at most $\rho_T$; for fixed $\mu$ the blocks of the inverse tend to $0$, so it is compact. The resolvent of $D_T$ is holomorphic on $N_\eta \supset \overline\Omega$, so its contour integral over $\Gamma$ vanishes, and the Riesz projection of $\hat D$ for $\Omega$ is the diagonal projection onto the window coordinates with $\lambda_j \in \Omega$.
\end{proof}

\begin{lemma}[small gain, Schur-complement form]\label{lem:smallgain}
Under the hypotheses of Lemma~\ref{lem:compare}, if $\theta_T < 1$, $\hat b < \infty$ and
\begin{equation}\label{eq:SC}
  f\!m_j + \frac{\hat b\,\rho_T\,r_j}{1 - \theta_T} < \operatorname{dist}_j \qquad\text{for every window column } j, \tag{SC}
\end{equation}
then $I - s\hat E(\mu - \hat D)^{-1}$ is invertible for every $(s, \mu) \in [0,1] \times \Gamma$.
\end{lemma}

\begin{proof}
Fix $(s, \mu)$. The tail block of $\mu - \hat D - s\hat E$ is $(I - s\hat E_{TT}(\mu - D_T)^{-1})(\mu - D_T)$, invertible with $\|(\mu - D_T - s\hat E_{TT})^{-1}\| \le \rho_T/(1 - \theta_T)$, since $\|\hat E_{TT}(\mu - D_T)^{-1}\| \le \theta_c\rho_T$ (the column sums of $\hat E_{TT}$ are at most $\sum_{n\ne0}\|A_n\| + \|A_0 - A_0^c\|$). By the block factorization (valid because the window block is finite-dimensional and both triangular factors are bijections of the domain), $\mu - \hat D - s\hat E$ is invertible if and only if the Schur complement
\[ \mathrm{Sc} = \mu - \Lambda - sR_W - s^2\hat E_{WT}(\mu - D_T - s\hat E_{TT})^{-1}\hat E_{TW} = (I - C_c(\mu - \Lambda)^{-1})(\mu - \Lambda) \]
is invertible, with $C_c = sR_W + s^2\hat E_{WT}(\cdots)^{-1}\hat E_{TW}$. Column $j$ of $\hat E_{TW}$ has norm at most $r_j$ (its tail rows are $\sum_w A_{m-w}V_{(w,\cdot),j}$, and $\sum_{m\in\mathrm{Tl}}\|A_{m-w}\| \le t_w$), and $\|\hat E_{WT}u\| \le \sum_{m,k}|u_{m,k}|b_m \le \hat b\|u\|$. So column $j$ of $C_c$ has norm at most $f\!m_j + \hat b\rho_Tr_j/(1 - \theta_T)$, and $\|C_c(\mu - \Lambda)^{-1}\| \le \max_j(f\!m_j + \hat b\rho_Tr_j/(1 - \theta_T))/\operatorname{dist}_j < 1$. Then $\mu - \hat D - s\hat E$ is invertible, which is equivalent to the claim because $\mu - \hat D$ is.
\end{proof}

The condition \eqref{eq:SC} asks only for the product of the window-to-tail coupling $\hat b$ and the tail-to-window coupling $r_j$ of each column to be small, which suits the problem: eigenvectors of near-axis eigenvalues are small at the window edge. The programs also evaluate the stronger form in which the window and tail columns are bounded separately; it does not hold at any $N$ and is not used.

\begin{lemma}[coupling sums from coefficient bounds]\label{lem:coupling}
Suppose enclosures $[A_n]$ of $A_n$ are known for $|n| \le n_A$ and $\|A_n\|_{1\to1} \le s_1q_1^{|n|} + s_2q_2^{|n|}$ for every $n$ ($q_1, q_2 < 1$). Let $G_{\mathrm{tail}}(k) = 2\sum_i s_iq_i^k/(1 - q_i)$. Then $\sigma_{\mathrm{off}} \le \sum_{0<|n|\le n_A}\|[A_n]\| + G_{\mathrm{tail}}(n_A + 1)$; $t_w \le \sum_{|n|\le n_A,\,|w+n|>K_e}\|[A_n]\| + G_{\mathrm{tail}}(n_A + 1)$; $b_m$ for $K_e < |m| \le K_e + n_c$ is bounded from entrywise enclosures of $A_{w-m}$ (for $|w - m| > n_A$ the entrywise tail enclosures of Section~\ref{sec:coeffs}); and for $|m| > K_e + n_c$, $b_m \le \beta_{\max}G_{\mathrm{tail}}(|m| - K_e)/2$, a bound that decreases in $|m|$, where $\beta_{\max} = \max_{(w,l)}\sum_j|(V^{-1})_{j,(w,l)}|$.
\end{lemma}

\begin{proof}
Sums of upper bounds, with exact geometric sums. For the last bound,
\[
\sum_{(w,l)}\beta_{(w,l)}|(A_{w-m})_{l,k}| \le \beta_{\max}\sum_w\|A_{w-m}\|_{1\to1},
\]
and the indices $n = w - m$ have one sign with $|n| \ge |m| - K_e$.
\end{proof}

\begin{lemma}[the window without inverting $V$]\label{lem:Vinv}
Let $V_i$ be an exact floating-point approximate inverse of $V$, $C = I - V_iV$, and suppose $q_C := \|C\| < 1$. Then $V$ is invertible, $V^{-1} = (I - C)^{-1}V_i$ with $\|(I - C)^{-1}\| \le 1/(1 - q_C)$, and with $\Xi := V_iH_{WW}V - \Lambda$,
\[ f\!m_j \le \frac{\|\Xi e_j\| + |\lambda_j|\,\|Ce_j\|}{1 - q_C}, \qquad \sum_j|(V^{-1})_{j,(w,l)}| \le \frac{\|V_ie_{(w,l)}\|}{1 - q_C}. \]
\end{lemma}

\begin{proof}
$V_iV = I - C$ is invertible by the Neumann series, so $V$ is invertible and $V^{-1} = (V_iV)^{-1}V_i$. Since $V_iV\Lambda = \Lambda - C\Lambda$, $V_i(H_{WW}V - V\Lambda) = \Xi + C\Lambda$ exactly, so $R_W = V^{-1}(H_{WW}V - V\Lambda) = (I - C)^{-1}(\Xi + C\Lambda)$, and $C\Lambda e_j = \lambda_jCe_j$.
\end{proof}

\begin{thmx}[the certificate]\label{thm:cert}
Suppose the following are checked in ball arithmetic:
\begin{enumerate}
\item[(C0)] $\omega \in [\omega_{\mathrm{lo}}, \omega_{\mathrm{hi}}]$, $\omega_{\mathrm{lo}} > 0$; entrywise enclosures $[A_n]$ of the true $A_n$ for every $n$ that enters (those of Section~\ref{sec:coeffs}, for $|n| \le n_A$ and beyond) and the tail form of Lemma~\ref{lem:coupling}; balls containing $d_m$;
\item[(C1)] $\delta > 0$; $R_0 > \alpha^{\mathrm{up}} \ge \alpha$; $a < 0 < b$, $b - a \ge \omega_{\mathrm{hi}}N$, $b < \omega_{\mathrm{lo}}N$, $-a < \omega_{\mathrm{lo}}N$;
\item[(C2)] the hypothesis of Lemma~\ref{lem:tail} for $r = 0, \dots, \lfloor N/2\rfloor$;
\item[(C3)] $V$ invertible and $f\!m_j$ bounded from a ball matrix containing the true $H_{WW}$ (with the ball $[\omega_{\mathrm{lo}}, \omega_{\mathrm{hi}}]$ on its diagonal); $\operatorname{dist}_j > 0$ for every $j$;
\item[(C4)] $\theta_T < 1$ and \eqref{eq:SC};
\item[(C5)] exactly one $\lambda_j$ lies in $\Omega$.
\end{enumerate}
Then \eqref{eq:G} holds.
\end{thmx}

\begin{proof}
By Lemmas~\ref{lem:compare} and~\ref{lem:smallgain}, Lemma~\ref{lem:homotopy} applies to $\hat D$ and $\hat E$, so $n(\mathrm{Scal}^{-1}H_0\mathrm{Scal}, \Omega) = n(\hat D, \Omega) = 1$ and $\Gamma$ lies in the resolvent set of $H_0$; the similarity by $\mathrm{Scal}$ gives $n(H_0, \Omega) = 1$. By (C1), $0 \in \Omega$, and $0$ is an eigenvalue (Lemma~\ref{lem:trivial}(a)), so it is the only eigenvalue in $\Omega$ and $m(0; H_0) = 1$. Let $\mu$ be an eigenvalue with $\re\mu \ge -\delta$. By Lemma~\ref{lem:H0}, $\re\mu \le \alpha < R_0$. Since $b - a \ge \omega_{\mathrm{hi}}N \ge \omega N$, some $\mu' = \mu + i\omega Nk$ has $a \le \im\mu' < a + \omega N \le b$, and $\mu'$ is an eigenvalue of the same multiplicity (Corollary~\ref{cor:strip}(i)). As $\Gamma$ is in the resolvent set, $\mu' \in \Omega$, so $\mu' = 0$, $\mu \in i\omega N\Z$ and $m(\mu; H_0) = 1$.
\end{proof}

The conditions $b < \omega_{\mathrm{lo}}N$ and $-a < \omega_{\mathrm{lo}}N$ are not needed for soundness (if $\Omega$ contained $\pm i\omega N$, the certificate could not pass, since by Theorem~\ref{thm:cert}'s argument $\Omega$ would contain at least two eigenvalues of $H_0$); they are needed for a correct spectrum to pass.

\subsection{From the existence ball to the coefficients}\label{sec:coeffs}

Theorem~\ref{thm:cert} needs enclosures of the coefficients $A_n$ of $Df(\phi_*(\theta))$ for the true profile, which is known only to lie in the ball of Theorem~\ref{thm:ring}(a). Let $J_n$ be the coefficients of $Df(\phi_{\bar a}(\theta))$ for the centre, enclosed as in Section~\ref{sec:fe} with the entrywise strip bound $|J_{n,kl}| \le S_{J,kl}e^{-\rho|n|}$, and let $M_k$ and $R_j$ be as in Section~\ref{sec:Z2}.

\begin{lemma}\label{lem:coeffs}
Let $r = r_{\mathrm{ex}}$, $\rho_e = \min(1/4, \rho_2) = 1/4$, $t_j = r$, assume $t_j < R_j$ for every $j$, and let
\[ \varepsilon_{kl} = \frac{M_k}{R_l}\Big[\Big(1 - \frac{t_l}{R_l}\Big)^{-1}\prod_j\Big(1 - \frac{t_j}{R_j}\Big)^{-1} - 1\Big]. \]
Then $|A_{kl}(\theta) - J_{kl}(\theta)| \le \varepsilon_{kl}$ for $|\im\theta| \le \rho_e$, and $|(A_n - J_n)_{kl}| \le \varepsilon_{kl}e^{-\rho_e|n|}$ for every $n$.
\end{lemma}

\begin{proof}
For $|\im\theta| \le 1/4$, $|\phi_{*,j}(\theta) - \phi_{\bar a,j}(\theta)| \le \|a_{*,j} - \bar a_j\|_\nu \le t_j$, so $w = \phi_* - \phi_{\bar a}$ lies in the polydisc of radii $t$. On the polydisc of radii $R$, $f_k(\phi_{\bar a}(\theta) + w) = \sum_\alpha c_\alpha w^\alpha$ with $|c_\alpha| \le M_kR^{-\alpha}$, so $\partial_lf_k(\phi_{\bar a} + w) - \partial_lf_k(\phi_{\bar a}) = \sum_{\alpha_l\ge1,|\alpha|\ge2}\alpha_lc_\alpha w^{\alpha - e_l}$ is bounded by the same sum for the majorant $M_k\prod_j(1 - t_j/R_j)^{-1}$, which is $\varepsilon_{kl}$. $A - J$ is holomorphic and $2\pi$-periodic on $|\im\theta| < 1/4$; Lemma~\ref{lem:cauchy} on every strip $|\im\theta| \le \rho''$, $\rho'' < 1/4$, and the limit $\rho'' \to 1/4$ give the coefficient bound.
\end{proof}

In the cell coordinates $S$ every bound transforms exactly ($(S^{-1}MS)_{kl} = M_{kl}2^{k_l - k_k}$), and from here on $[J_n]$, $S_J$ and $\varepsilon$ denote the transformed bounds. The radius must be $r_{\mathrm{ex}}$, not $r_{\mathrm{un}}$: at $N = 8$ the record gives $\|\varepsilon\|_{1\to1} \le 1.2 \times 10^{-12}$ in the coordinates $S$ with $r_{\mathrm{ex}}$; since $\varepsilon$ grows linearly in $r$ for small $r$, the radius $r_{\mathrm{un}} = 10^{-12}$ would give about $7 \times 10^3$ in the same coordinates, far above the near-axis margins $\operatorname{dist}_j \approx 1.3 \times 10^{-6}$. So the program uses the entrywise enclosures $[A_n] = [J_n] + \operatorname{ball}(0, \varepsilon e^{-\rho_e|n|})$ for $|n| \le n_A = 80$ and $[A_n] = \operatorname{ball}(0, S_Je^{-\rho|n|} + \varepsilon e^{-\rho_e|n|})$ for $|n| > 80$ (entrywise; this holds for every $n$ by Lemma~\ref{lem:cauchy} applied to $J$ on the strip $|\im\theta| \le \rho$ and by Lemma~\ref{lem:coeffs}), which is needed in the window ($|n| \le 2K_e$, up to $96$ at $N = 64$) and for $b_m$ ($|n| \le 2K_e + n_c$, up to $120$), together with the tail form $\|A_n\|_{1\to1} \le \|S_J\|e^{-\rho|n|} + \|\varepsilon\|e^{-\rho_e|n|}$ for every $n$, which is the form of Lemma~\ref{lem:coupling}.

\subsection{Conclusion for the wave}\label{sec:conclusion}

\begin{thmx}\label{thm:stab}
Assume the existence statement of Lemma~\ref{lem:real} and the certificate of Theorem~\ref{thm:cert} for some $\delta > 0$, with inputs (C0) built as in Section~\ref{sec:coeffs}. Then:
\begin{enumerate}
\item[(i)] $Y(T)$ has the eigenvalue $1$ algebraically simple, and the other $18N - 1$ Floquet multipliers satisfy $|\rho| < e^{-\delta T} \le e^{-\delta T_{\mathrm{lo}}}$, $T_{\mathrm{lo}} = 2\pi/\omega_{\mathrm{hi}}$;
\item[(ii)] $M_\tau$ has the eigenvalue $1$ algebraically simple and all others of modulus less than $e^{-\delta\tau}$, and so has every section map of Corollary~\ref{cor:section}, without the eigenvalue $1$;
\item[(iii)] the orbit is locally exponentially orbitally stable with asymptotic phase, as stated in Theorem~\ref{thm:ring}(c).
\end{enumerate}
\end{thmx}

\begin{proof}
(i) By Corollary~\ref{cor:strip}(iii) with $a = -\omega/2$, the multipliers are $e^{\mu T}$, $\mu \in \spec H_0 \cap S_a$, with multiplicities. By \eqref{eq:G}, $\mu = 0$ is simple and is the only element of $i\omega\Z \cap S_a$ in the spectrum, and every other $\mu$ has $\re\mu < -\delta$, so $|e^{\mu T}| < e^{-\delta T}$. Lemma~\ref{lem:trivial}(b) gives simplicity of the multiplier $1$. (ii) Corollaries~\ref{cor:strip}(ii) and~\ref{cor:section}.

(iii) $F$ is analytic near $O = \{x(t)\}$. Let $\Sigma = \{z : \langle F(x^*), z - x^*\rangle = 0\}$. Since $\frac{d}{dt}\langle F(x^*), \varphi_t(z) - x^*\rangle$ at $(T, x^*)$ equals $|F(x^*)|^2 > 0$, the implicit function theorem gives a neighbourhood $U$ of $x^*$ and a $C^1$ function $t_S$ on $U$ with $t_S(x^*) = T$ and $\varphi_{t_S(z)}(z) \in \Sigma$. The Poincar\'e map $P(z) = \varphi_{t_S(z)}(z)$ on $\Sigma \cap U$ has $P(x^*) = x^*$ and, as in Corollary~\ref{cor:section} with $Q$ replaced by $I$ and $\tau$ by $T$, $DP(x^*)$ is the compression of $Y(T)$, with spectrum that of $Y(T)$ minus one copy of $1$; by (i) its spectral radius $\varrho$ is less than $e^{-\delta T}$. Let $\delta' < \delta$, $\kappa = e^{-\delta'T}$, and choose $\kappa'$ with $\varrho < \kappa' < \kappa$. With a Jordan form $J = U_0^{-1}DP(x^*)U_0$ and $D_\epsilon = \operatorname{diag}(1, \epsilon_J, \epsilon_J^2, \dots)$, the norm $|v|_* = |D_\epsilon^{-1}U_0^{-1}v|_\infty$ with $0 < \epsilon_J \le \kappa' - \varrho$ gives $\|DP(x^*)\|_* \le \kappa'$. By continuity there is $\epsilon_1 > 0$ with $\|DP(z)\|_* \le \kappa$ on the ball $B_c = \{z \in \Sigma : |z - x^*|_* \le \epsilon_1\}$, which is convex; the mean value inequality gives $|P(z) - x^*|_* \le \kappa|z - x^*|_*$ there, so $P$ maps $B_c$ into itself and $|P^k(z) - x^*|_* \le \kappa^k|z - x^*|_*$.

Let $x^0$ be given and choose $s \in [0, T)$ with $|x^0 - x(s)| = \operatorname{dist}(x^0, O)$ (it exists because $O$ is compact), and $z_0 = \varphi_{T-s}(x^0)$; by Gronwall $|z_0 - x^*| \le e^{\mathrm{Lip}\,T}|x^0 - x(s)|$, so the solution meets $\Sigma$ at time $t_1 = T - s + t_S(z_0) \le 3T$ at $z_1 = \varphi_{t_S(z_0)}(z_0)$ with $|z_1 - x^*| \le C_1\operatorname{dist}(x^0, O)$, and we take $\epsilon_0$ so small that $z_1 \in B_c$. Let $z_k = P^{k-1}(z_1)$ and $t_{k+1} = t_k + t_S(z_k)$. Since $|t_S(z_k) - T| \le C_2\kappa^k|z_1 - x^*|$, the limit $\sigma_\infty = \lim(t_k - kT)$ exists with $|t_k - kT - \sigma_\infty| \le C_3\kappa^k|z_1 - x^*|$. Put $\sigma = -\sigma_\infty$. For $t \in [t_k, t_{k+1}]$, $\varphi_t(x^0) = \varphi_{t-t_k}(z_k)$ and $x(t + \sigma) = \varphi_{t-t_k}(x(t_k + \sigma))$, where $x(t_k + \sigma) = x(t_k - kT - \sigma_\infty)$ lies within $\sup|F|\,C_3\kappa^k|z_1 - x^*|$ of $x^*$ and $|z_k - x^*| \le C_4\kappa^k|z_1 - x^*|$; Gronwall over an interval of length at most $2T$ gives $|\varphi_t(x^0) - x(t + \sigma)| \le C_5\kappa^k|z_1 - x^*|$, and $t \le (k+1)T + C_6$ gives $\kappa^k \le C_7e^{-\delta't}$. On $[0, t_1]$, Gronwall gives $|\varphi_t(x^0) - x(t + s)| \le e^{\mathrm{Lip}\,t}|x^0 - x(s)|$, and since $|\sigma_\infty - (t_1 - T)| \le \sum_{k\ge1}|t_S(z_k) - T| \le C'|z_1 - x^*|$ and $t_1 - T = -s + t_S(z_0)$ with $|t_S(z_0) - T| \le C|z_0 - x^*|$, $\sigma = s - T + O(\operatorname{dist}(x^0, O))$, so by $T$-periodicity $|x(t + \sigma) - x(t + s)| \le C''\operatorname{dist}(x^0, O)$. Together, on $[0, t_1]$ the difference is at most $C_8e^{3\delta'T}e^{-\delta't}\operatorname{dist}(x^0, O)$.
\end{proof}

This is the classical Andronov--Witt statement, proved here from the Poincar\'e map so that no textbook hypothesis needs matching. If $T$ were not the minimal period of the ring solution, (i) to (iii) would still hold with respect to $T$; Lemma~\ref{lem:real} shows that it is minimal.

\subsection{The computations}\label{sec:stabcomp}

The program \file{code/fourier/stability.py} implements Theorem~\ref{thm:cert} with Lemma~\ref{lem:tail} for (C2), Lemma~\ref{lem:smallgain} for (C4) and Lemma~\ref{lem:Vinv} for (C3), at 128 bits everywhere. It reads $\omega_{\mathrm{lo}}$, $\omega_{\mathrm{hi}}$, $r_{\mathrm{ex}}$ and the settings from the Stage~E record (it refuses to run if any source hash in that record differs from the current file, and it stores the SHA-256 of the record it read), and recomputes $[J_n]$, $S_J$ and $M_k$ by the same calls as the existence program; as a consistency check their suprema and the digests of the trigonometric polynomials must equal the values in the Stage~E record. It asserts $t_j < R_j$; rounds $\delta$ up to a dyadic with denominator $2^{60}$; takes $R_0 = 32$, the least power of two above $\alpha^{\mathrm{up}}$, which lies between $17.84$ and $18.11$ for the five values of $N$; takes $b = -a$ an exact double near $\bar\omega(N/2 + 1/4)$, which puts both horizontal edges a quarter of $\omega$ from the rows $\im\mu \approx \omega k$ where the near-axis eigenvalues lie; takes $K_e = \lfloor N/2\rfloor + 16$, $n_c = 24$, $\eta = 2^{-20}$; uses Lemma~\ref{lem:tail} for the tail (the program calls this ``route A''; an alternative ``route B'', a direct cover of the bounded part described in \file{code/fourier/LEMMAS-stability.md}, is not used); chooses the exponents of $S$ in floating point to make the tail bound small; computes $V$ and $\Lambda$ by LAPACK on the midpoint of the window block in the coordinates $S$, single-threaded; and decides every inequality of (C1) to (C5) in Arb. The count (C5) compares the exact doubles $\lambda_j$ with the exact dyadic edges $-\delta$, $R_0$, $a$, $b$. Floating point only chooses exact numbers, which decides whether the inequalities hold, never whether a passed inequality is true. Before any conclusion the program also checks that the residual of the trivial eigenvector $P^*$ under the window block is small (about $10^{-20}$), which would expose a dropped coefficient or a wrong sign convention.

\begin{proof}[Proof of Theorem~\ref{thm:ring}(c) and Theorem~\ref{thm:cell}(ii)]
For each $N \in \{1, 8, 16, 32, 64\}$ the records \file{data/fourier-stability-N*.json} show (C0) to (C5) passed, with Lemma~\ref{lem:tail} for the tail and the form (SC): for example at $N = 8$, $\rho_T = 3.578$, $\theta_c = 0.0787$, $\theta_T = 0.2817$, $\hat b = 0.856$, $q_C = 2.6 \times 10^{-14}$, the smallest $\operatorname{dist}_j$ is $1.32 \times 10^{-6}$, and exactly one $\lambda_j$ (at about $3 \times 10^{-16} - 4 \times 10^{-15}i$) lies in $\Omega$; Table~\ref{tab:stability} lists the other $N$. Theorem~\ref{thm:stab} gives the statements, with the multiplier bounds $e^{-\delta T_{\mathrm{lo}}}$ and $e^{-\delta T_{\mathrm{lo}}/N}$ rounded upward from Arb. For $N = 1$, $Y(T) = M_\tau$ and $Q = I$.
\end{proof}

\section{The cell in the time domain}\label{sec:capd}

The proof of Theorem~\ref{thm:cell}(i) uses CAPD \cite{capd2021}. Work in the scaled variables $z$ of Section~\ref{sec:cell}, and let $\mathcal S = \{z_V = s/\sigma_V\}$ (that is $V = s$; $s/\sigma_V = 4s$ is an exact double) and $g$ the first-return map to $\mathcal S$, crossed upward. The frame file \file{code/candidates/cell_frameF.txt}, whose SHA-256 the record \file{data/cell-gks0.0275.json} stores, gives a point $\hat x$ with $\hat x_V = s/\sigma_V$ exactly (the verifier refuses the file otherwise), an invertible $17 \times 17$ matrix $A_t$, and a partition of the 17 free coordinates into 14 blocks of size 1 or 2 with radii $\rho_b$ between $1.13 \times 10^{-13}$ and $8.36 \times 10^{-11}$ (listed in the record). With $A = \begin{pmatrix}1 & 0\\ 0 & A_t\end{pmatrix}$ the map $y \mapsto \hat x + A(0, y)$, $y \in \R^{17}$, is an affine parametrization of $\mathcal S$, since the first component of $A(0, y)$ is $0$. Let $\|y\| = \max_b\|y_b\|_2/\rho_b$ and
\[ B = \{\hat x + A(0, y) : \|y\| \le 1\}, \qquad G(y) = A_t^{-1}\big(g(\hat x + A(0, y)) - \hat x\big)_{\mathrm{free}}. \]
$\hat x$ and $A_t$ are proposed by the untrusted programs \file{code/proofs/orbit_newton.cpp} and \file{code/proofs/frame.py} (eigenvectors of the numerical $Dg$ for the slow directions and an orthonormal Schur basis for the fast ones). The program \file{code/proofs/verify.cpp} encloses $A_t^{-1}$ rigorously, $G(0)$ by a 128-bit MPFR interval run (order 30) and $DG$ over the whole box $\prod_b[-\rho_b, \rho_b] \supset \{\|y\| \le 1\}$ by a $C^1$ Lohner run in double intervals (order 20), and checks:
\begin{itemize}
\item $q = \max_b\sum_c\|M_{bc}\|\rho_c/\rho_b < 1$, where $M$ encloses $DG$ over the box (exact spectral bounds on diagonal $2 \times 2$ blocks, Frobenius bounds elsewhere);
\item $\|G(0)_b\|/\rho_b + (\text{row sum of block } b) < 1$ for every block $b$;
\item every bound entering the decision is finite, and the section time of the centre run lies inside the section time of the box run, so both resolve the same crossing; CAPD's crossing step checks that $dV/dt$ does not vanish on the enclosure of the crossing point, so the crossing is transversal and $g$ is $C^1$ on $B$.
\end{itemize}
Then $G$ maps $\{\|y\| \le 1\}$ into itself and is a $q$-contraction there in the norm $\|\cdot\|$ (mean value inequality on a convex set), so $g$ has a unique fixed point $x_c$ in $B$ and every eigenvalue of $Dg(x_c)$ has modulus at most $q$. By Corollary~\ref{cor:section} with $N = 1$ ($Q = I$, $\tau = T$), $\det(\lambda I - Y(T)) = (\lambda - 1)\det(\lambda I - Dg(x_c))$, so these eigenvalues are the 17 nontrivial Floquet multipliers. Local orbital asymptotic stability, with asymptotic phase and rate $e^{-\delta' t}$ for every $\delta'$ with $e^{-\delta' T} > q$, follows from the proof of Theorem~\ref{thm:stab}(iii) with $\Sigma$ replaced by $\mathcal S$ near $x_c$, $P$ by $g$ (whose return time is $C^1$ by transversality), and the ball $B_c$ by $B$ with the norm $\|\cdot\|$, in which $g$ is already a $q$-contraction. The period is the return time: since $g(x_c) = x_c$, the first return time $t_1$ of $x_c$ is a period, and since the orbit is back at $x_c \in \mathcal S$ at its minimal period, $t_1$ is the minimal period. The fixed point lies in the relative interior of $B$ in $\mathcal S$, as the proof of Theorem~\ref{thm:stab}(iii) needs: by the second check, $G$ maps $\{\|y\| \le 1\}$ into $\{\|y\| \le 0.99865\}$, and $x_c$ corresponds to a fixed point of $G$. The record gives $q \le 0.99864150686405151$ (stored as the double \texttt{0x1.ff4df088d18aap-1}), $\max_b(\|G(0)_b\|/\rho_b + \text{row sum}) \le 0.99864258833469344$, and the return time $[53.5855190480139, 53.585519630722438]$~ms; the run took $732$~s.

\begin{lemma}[the two cell certificates enclose the same orbit]\label{lem:link}
Let $\phi_*$ be the profile of Theorem~\ref{thm:cell}(ii) and $p = \Sigma^{-1}\phi_*(0)$. Then $p \in B$, and the orbit $\{\Sigma^{-1}\phi_*(\theta)\}$ of Theorem~\ref{thm:cell}(ii), in scaled variables, is the orbit of the fixed point $x_c$ of Theorem~\ref{thm:cell}(i).
\end{lemma}

\begin{proof}
$p_V = s/\sigma_V$ by the phase condition, so $p \in \mathcal S$. For each free component $k$, $|p_k - \bar p_k| \le \sum_m|a_{*,k,m} - \bar a_{k,m}| \le \|a_{*,k} - \bar a_k\|_\nu \le r_{\mathrm{ex}}$, where $\bar p = \sum_{|m|\le K}\bar a_m$ is an exact rational vector (the program checks that $\operatorname{Im}\bar a_0 = 0$, so $\bar p = \bar a_0 + 2\re\sum_{m=1}^K\bar a_m$). The program \file{code/fourier/link_cell.py} reads $\hat x$, $A_t$ and the radii as the doubles the verifier reads, inverts $A_t$ exactly in rational arithmetic, computes $y^c = A_t^{-1}(\bar p - \hat x)_{\mathrm{free}}$ exactly and the bound $|y_k - y^c_k| \le r_{\mathrm{ex}}\sum_j|(A_t^{-1})_{kj}|$ for $y = A_t^{-1}(p - \hat x)_{\mathrm{free}}$, and checks, with upper bounds of square roots in integer arithmetic, $\|y^c_b\|_2 + \|(y - y^c)_b\|_2 < \rho_b$ for every block. All 14 checks pass, the largest ratio being below $0.14$ (\file{data/link_cell.txt}); a control with every radius divided by $10$ fails. So $p \in B$. Since $g(B) \subset B$ and $g$ is a contraction on $B$, $g^k(p) \in B$ and $g^k(p) \to x_c$. Each $g^k(p)$ is a point of the flow line through $p$ of the CAPD field; by the assumptions of Theorem~\ref{thm:cell}(i) and (ii) taken together, the CAPD and Arb fields are the same function, so that flow line is the closed set $\{\Sigma^{-1}\phi_*(\theta)\}$. So $x_c$ lies on it, and the orbit of $x_c$ is that set. Neither the crossing direction of $p$ nor $x_c = p$ is used: CAPD defines $g$ on all of $B$, and only that $g^k(p)$ lies on the flow line of $p$ is needed.
\end{proof}

\paragraph{The domain and the CAPD patch.} CAPD's interval wrapper throws when a divisor contains $0$, a logarithm argument is not positive or a square-root argument is negative, and the run then fails; the multiprecision logarithm does not check its argument, so every $\log a$ is computed as $2\log\sqrt a$ (the same function for $a > 0$), and the multiprecision square root refuses a negative argument. The source's $(1 - u)$ is computed as $1/(1 + e^{5(V+40)})$, the same function, because near $V = 0$, $u$ is within $10^{-80}$ of $1$. In \texttt{PoincareMap::crossSectionInOneStep} of CAPD 6.1.0 the crossing-point bound is trimmed with the endpoints of the previous Newton window while monotonicity was checked on the current one, which is unsound if the loop stops at its iteration cap without converging. The run used a header-only local patch that takes the endpoints of the current window (\file{code/proofs/capd-6.1.0-genchase.patch}); the record stores the SHA-256 of the patched header. An earlier unpatched run gave the same bounds.

\section{Numerical observations}\label{sec:numerics}

Nothing in this section is proved, and none of it is used in a proof.

\begin{itemize}
\item \emph{Leading exponents.} The floating-point eigenvalues of the window blocks closest to the imaginary axis, apart from $0$, have real parts $-4.71382 \times 10^{-5}$ ($N = 1$), $-6.32096 \times 10^{-6}$ ($N = 8$), $-8.57417 \times 10^{-6}$ ($16$), $-9.18618 \times 10^{-6}$ ($32$) and $-9.34235 \times 10^{-6}$ per ms ($64$), with imaginary parts near $0$ ($N = 1$) and near $\pm\omega$ (the rings). The certified $\delta = 5 \times 10^{-6}$ for the rings is below all of them.
\item \emph{Sharpness at $N = 8$.} In runs of \file{code/fourier/stability.py} made during the in-project reviews and in its tests, not kept as records, the certificate passes at $N = 8$ with $\delta = 6.32095 \times 10^{-6}$ (near-axis (SC) ratio $0.398$) and fails at $\delta = 6.321 \times 10^{-6}$ with a count of $3$ in (C5). So the certificate is sharp to about $10^{-6}$ relative against the floating-point exponent; we do not state the larger $\delta$ as a theorem.
\item \emph{Negative controls.} In the same runs the certificate failed, at the step that is supposed to fail, for $\delta = 7 \times 10^{-6}$ at $N = 8$ ((C5)), for reversed coupling $c \to -c$ ((C5), count $9$), for a window that is too small (Lemma~\ref{lem:tail}), for $S = I$ (Lemma~\ref{lem:tail} or $\theta_T < 1$ fails), for a coefficient $A_{\pm1}$ dropped from the proof data only ((SC) fails in 127 columns), and for floating-point data that place the leading pair at $-7.5 \times 10^{-6}$ while the true data do not ((SC) fails).
\item \emph{Hopf point.} An ordinary finite-difference computation (\file{code/numerics/hopf_and_orbit.py}, \file{data/numerics-hopf-orbit.json}) reproduces Erhardt's Hopf point to $5 \times 10^{-7}$ relative ($0.0279078439$ against $0.0279078589$) and the Hopf frequency to $3 \times 10^{-7}$ relative. In that computation the voltage of the cell orbit stays between about $0.14$ and $0.26$~mV: a small depolarized oscillation, not an action potential.
\item \emph{Ring periods.} The certified periods increase with $N$ by $2.45 \times 10^{-3}$, $9.8 \times 10^{-5}$, $2.5 \times 10^{-5}$ and $6.3 \times 10^{-6}$~ms from $N = 1$ to $8$, $8$ to $16$, $16$ to $32$ and $32$ to $64$. The ratios of successive differences, about $3.9$ and $4.0$ for the last two, look like second-order convergence in $1/N$; that is an observation about four numbers, not a continuum result.
\end{itemize}

\section{Reproducibility and trust base}\label{sec:reproduce}

The programs are in \file{code/}, in the layout of their canonical location \file{research/cardiac-cycle-certificates/} of the GENChase repository, whose records hash them; the records in \file{data/} are byte-identical copies, and \file{code/fourier/check_records.py} rechecks all 90 stored hashes (sources, centres and the Stage~E records read by Stage~S) in a few seconds. \file{code/run_all.sh} stages the programs and records in a scratch folder, runs that check and the link of Lemma~\ref{lem:link} (\file{code/fourier/link_cell.py}, standard library only), and reruns Stage~E and Stage~S for chosen $N$.

\paragraph{Fourier route.} Trusted: Arb and FLINT through python-flint 0.9.0 (FLINT 3.6.0; the wheel's SHA-256 is pinned and checked by \file{code/fourier/test_arbmodel.py}), CPython, the programs \file{code/fourier/arbmodel.py}, the generated \file{code/fourier/tp06_18d_arb.py}, \file{fourier_eval.py}, \file{existence.py}, \file{stability.py}, and the reference translation \file{code/model/tp06_18d.py} with the scales. The generator rewrites every decimal literal of the reference translation as an exact rational by Python's tokenizer, and a freshness check compares the generated file byte for byte. numpy and LAPACK only propose centres, eigenvectors and weights. The tests compare the Arb field with an mpmath evaluation and with the CAPD field at random points, test the strip cover, the aliasing bound and the centred form against closed forms, and contain the negative controls of Section~\ref{sec:numerics}, independent floating-point estimates that every certified bound must dominate, and perturbation controls within the uniqueness ball.

\paragraph{CAPD route.} Trusted: CAPD 6.1.0 at commit \texttt{03dc5628203334b214bb7d9fd63788a175521005} with the local patch above, filib and MPFR, g++ 13.3 with \texttt{-frounding-math}, the processor's directed rounding, and \file{code/proofs/verify.cpp} with \file{code/model/tp06_capd.hpp} and \file{setup.hpp}, in which every non-integer decimal enters as an interval enclosing that exact decimal. The driver \file{code/proofs/certify.py} checks that the frame file's header names the intended claim and records the hashes of the frame, the binary and the sources.

\paragraph{What is certified, and reruns.} The theorems rest on the stored records in \file{data/}: the Stage~S records hash the stored Stage~E records they read, so the stored $r_{\mathrm{ex}}$ is the one consumed. On 2026-10-01 the existence and stability programs for $N = 1$ and $N = 8$ were rerun from the copies in \file{code/} with python-flint 0.9.0 (\file{notes/rerun-2026-10-01.md}): the period and frequency enclosures, $r_{\mathrm{un}}$ and every stability bound of Table~\ref{tab:stability} agree exactly with the records, while the binary values of $Y_0$, $Z_1$, $Z_2$ and $r_{\mathrm{ex}}$ differ from the records in late digits (for $Y_0$ beyond the 25 printed digits, for the others in about the thirteenth significant digit). The likely cause is that \file{existence.py}, unlike \file{stability.py}, does not pin the number of BLAS threads, so the untrusted floating-point inverse $A_{\mathrm{fin}}$ need not be bit-reproducible; any exact $A_{\mathrm{fin}}$ is admissible, and only the size of the bounds depends on it. A fix is queued in the project. An earlier in-project reading reported a bit-identical Stage~E rerun on the machine where the records were made. The rings $N = 16$, $32$, $64$ and the CAPD certificate were not rerun from these copies.

The run times quoted are those of the records: about two minutes for each existence proof, and $26$, $48$, $80$, $184$ and $479$~s for the stability proofs at $N = 1, 8, 16, 32, 64$ (peak memory $3.6$~GB at $N = 64$).

\section{Limitations and what has been checked}\label{sec:limits}

\begin{itemize}
\item \emph{Scope.} The theorems are about one modified cell model at one parameter value, $G_{Ks} = 0.0275$, and about four ring sizes with one coupling law. Nothing is claimed for the published 19-state TP06 cell, for action potentials or reentry, for a continuum cable, for tissue, for other $N$ or uniformly in $N$. The orbits are small depolarized oscillations (Section~\ref{sec:numerics}).
\item \emph{The Hopf branch.} That the cell orbit lies on the branch of periodic orbits born at Erhardt's Hopf point, and that the ring waves continue it, is supported by numerics only.
\item \emph{Future work: an interval in $G_{Ks}$.} A certified branch of these orbits on an interval of $G_{Ks}$ toward the Hopf point is in progress in the project (\file{research/cardiac-cycle-certificates/fourier/branch.py}); no result of it is used or claimed here.
\item \emph{Translations of the model.} Theorem~\ref{thm:cell}(i) assumes that the CAPD field (\file{code/model/tp06_capd.hpp}) evaluates Erhardt's \texttt{fun\_eval} with exact decimals, and (ii) and Theorem~\ref{thm:ring} assume the same of the Arb field (\file{code/fourier/tp06_18d_arb.py}); Lemma~\ref{lem:link} uses both assumptions together. The tests compare the two programs' enclosures at 104 points, which is evidence, not a proof that they define the same function.
\item \emph{Kato.} The facts of Section~\ref{sec:hill} are cited from Kato's book by chapter and section; the book has not been checked against a copy for the exact statements or their numbering.
\item \emph{How far the background sources were read.} None of them is a source of a proof step; only Erhardt's model source and \cite{erhardt2025} (read in full) enter the setting. Read in part: Gameiro and Lessard \cite{gl2017} (Sections 1, 7 and 8 of the authors' version), Bayer and Leine \cite{bl2026} (Sections 1 to 4, 6 and 7 of arXiv:2503.21318v2), Ermentrout \cite{ermentrout1992} (pp.~1674--1677), Johnson and Zumbrun \cite{jz2012} (first page), Czechowski and Zgliczy\'nski \cite{cz2016} (Theorems 1.1 to 1.3). Known from abstracts: Paullet and Ermentrout \cite{pe1994}, Church and Lessard \cite{cl2022}, Arioli and Koch \cite{ak2020}, Kapela and Zgliczy\'nski \cite{kz2003}, Kuehn and Queirolo \cite{kq2022}, ten Tusscher and Panfilov \cite{tp2006}, Church et al.\ \cite{cdhlv2026}, Castelli and Lessard \cite{cl2013}, Rucklidge and Silber \cite{rs1998}, de Wolff \cite{dewolff2023}, the papers of Erhardt and Solem \cite{es2021,es2022} and of K\"ugler et al.\ \cite{keb2017,keb2018}. Known from its title and Czechowski and Zgliczy\'nski's account: Arioli and Koch \cite{ak2015}. Known from metadata or a table of contents: Golubitsky, Stewart and Schaeffer \cite{gss1988} (table of contents), Golubitsky and Stewart \cite{gs1985}, Figueras et al.\ \cite{fgld2017}, Courtemanche et al.\ \cite{cgk1993}, Pecora and Carroll \cite{pc1998}, Ashwin and Swift \cite{as1992}, Hoppensteadt and Izhikevich \cite{hi1997}. Known from excerpts of its abstract: Di Marco et al.\ \cite{dmfgkp2016}. The ledger and \file{research/cardiac-cycle-certificates/notes/} record each reading.
\item \emph{Searches.} The prior-article searches are logged in the GENChase research ledger (\file{RESEARCH.md}, entries of 2026-09-30 and 2026-10-01 on cardiac work) and in \file{research/cardiac-cycle-certificates/notes/}. The ledger asks that Paullet and Ermentrout \cite{pe1994}, Ashwin and Swift \cite{as1992} and Hoppensteadt and Izhikevich \cite{hi1997} (Theorem~9.2) be read in full, and the papers citing Erhardt \cite{erhardt2025} be checked for a continuation of the first Hopf branch, before submission; that has not been done.
\item \emph{Checks.} Each stage was read adversarially within the project by separate AI agent sessions: the Fourier evaluation and model code, the stability lemmas, the existence program and the stability program, followed by a second reading of the fixes, which found nothing unsound (\file{review/}). The CAPD verifier had a five-lens adversarial review within the project, which led to the patch above. A further in-project reading of this manuscript (\file{review/manuscript-reading-1-2026-10-01.md}) found no false theorem; its findings on rounding, wording, precision of the existence bounds, citations and exposition were addressed. A second reading of the revised draft (\file{review/manuscript-reading-2-2026-10-01.md}) found nothing unsound, confirmed the link of Lemma~\ref{lem:link} by rerunning it, and listed corrections that are made in this version (\file{review/fix-check-2026-10-01.md}). Three items of the program readings remain open: a few negligible terms (for example the strip tail of $S_J$ beyond $K'$) cannot be detected by any test; the source hashes are taken after the modules are imported, a small window; and the Kato numbering above.
\end{itemize}

\section{Discussion}\label{sec:discussion}

The method separates the dimension of the ring from the cost of the proof: existence costs the same for every $N$, and stability grows with the window $K_e \approx N/2 + 16$ that must contain the $N$ near-axis copies of the spectrum. The Hill-sector theorem turns a $18N$-dimensional Floquet problem into a question about one operator on $\ell^1(\Z; \C^{18})$, and the Riesz homotopy with the Schur-complement small gain answers it with a count, which is what a stability proof needs and an eigenpair enclosure cannot give. The margins are thin: at $N = 8$ the leading nontrivial exponent is about $-6.32 \times 10^{-6}$ per ms, so the multipliers are within about $3.4 \times 10^{-4}$ of the unit circle per period; the certificate still resolves them because the window is computed at 128 bits and the near-axis columns of (SC) have ratios below $10^{-6}$. Natural next steps are an interval of $G_{Ks}$ (in progress), larger $N$, and statements uniform in $N$.

\medskip
\noindent\textbf{Funding.} This research received no external funding.

\noindent\textbf{Use of AI.} This work was prepared with AI assistance. The author takes full responsibility for its content.

\begin{thebibliography}{99}

\bibitem{ak2015} G. Arioli and H. Koch, Existence and stability of traveling pulse solutions of the FitzHugh--Nagumo equation, \emph{Nonlinear Anal.} \textbf{113} (2015) 51--70; doi:10.1016/j.na.2014.09.023.

\bibitem{ak2020} G. Arioli and H. Koch, Traveling wave solutions for the FPU chain: a constructive approach, \emph{Nonlinearity} \textbf{33} (2020) 1705--1722; doi:10.1088/1361-6544/ab6a78.

\bibitem{as1992} P. Ashwin and J. W. Swift, The dynamics of $n$ weakly coupled identical oscillators, \emph{J. Nonlinear Sci.} \textbf{2} (1992) 69--108; doi:10.1007/BF02429852.

\bibitem{bl2026} F. Bayer and R. I. Leine, Explicit error bounds and guaranteed convergence of the Koopman--Hill projection stability method for linear time-periodic dynamics, \emph{J. Nonlinear Sci.} \textbf{36} (2026) 64; doi:10.1007/s00332-026-10287-3; arXiv:2503.21318v2, whose sections are cited.

\bibitem{cl2013} R. Castelli and J.-P. Lessard, Rigorous numerics in Floquet theory: computing stable and unstable bundles of periodic orbits, \emph{SIAM J. Appl. Dyn. Syst.} \textbf{12} (2013) 204--245; doi:10.1137/120873960.

\bibitem{cdhlv2026} K. E. M. Church, J.-Y. Dai, O. H\'enot, P. Lappicy and N. Vassena, Global continuation of stable periodic orbits in systems of competing predators, \emph{SIAM J. Appl. Dyn. Syst.} \textbf{25} (2026) 1697--1725; doi:10.1137/25M1748275; arXiv:2504.03058.

\bibitem{cl2022} K. E. M. Church and J.-P. Lessard, Rigorous verification of Hopf bifurcations in functional differential equations of mixed type, \emph{Physica D} \textbf{429} (2022) 133072; doi:10.1016/j.physd.2021.133072.

\bibitem{cgk1993} M. Courtemanche, L. Glass and J. P. Keener, Instabilities of a propagating pulse in a ring of excitable media, \emph{Phys. Rev. Lett.} \textbf{70} (1993) 2182--2185; doi:10.1103/PhysRevLett.70.2182.

\bibitem{cz2016} A. Czechowski and P. Zgliczy\'nski, Existence of periodic solutions of the FitzHugh--Nagumo equations for an explicit range of the small parameter, \emph{SIAM J. Appl. Dyn. Syst.} \textbf{15} (2016) 1615--1655; doi:10.1137/15M1007707; arXiv:1502.02451.

\bibitem{dewolff2023} B. de Wolff, Equivariant Pyragas control of discrete waves, \emph{SIAM J. Math. Anal.} \textbf{55} (2023) 6707--6739; doi:10.1137/22M1527143.

\bibitem{dmfgkp2016} M. Di Marco, M. Forti, B. M. Garay, M. Koller and L. Pancioni, Floquet multipliers of a metastable rotating wave in a Chua--Yang ring network, \emph{J. Math. Anal. Appl.} \textbf{434} (2016) 798--836; doi:10.1016/j.jmaa.2015.08.072.

\bibitem{erhardt2025} A. H. Erhardt, Cardiac dynamics of a human ventricular tissue model with a focus on early afterdepolarizations, \emph{Front. Phys.} \textbf{13} (2025) 1569121; doi:10.3389/fphy.2025.1569121.

\bibitem{es2021} A. H. Erhardt and S. Solem, On complex dynamics in a Purkinje and a ventricular cardiac cell model, \emph{Commun. Nonlinear Sci. Numer. Simul.} \textbf{93} (2021) 105511; doi:10.1016/j.cnsns.2020.105511.

\bibitem{es2022} A. H. Erhardt and S. Solem, Bifurcation analysis of a modified cardiac cell model, \emph{SIAM J. Appl. Dyn. Syst.} \textbf{21} (2022) 231--247; doi:10.1137/21M1425359; preprint arXiv:2105.05544.

\bibitem{ermentrout1992} G. B. Ermentrout, Stable periodic solutions to discrete and continuum arrays of weakly coupled nonlinear oscillators, \emph{SIAM J. Appl. Math.} \textbf{52} (1992) 1665--1687; doi:10.1137/0152096.

\bibitem{fgld2017} J.-L. Figueras, M. Gameiro, J.-P. Lessard and R. de la Llave, A framework for the numerical computation and a posteriori verification of invariant objects of evolution equations, \emph{SIAM J. Appl. Dyn. Syst.} \textbf{16} (2017) 1070--1088; doi:10.1137/16M1073777.

\bibitem{gl2017} M. Gameiro and J.-P. Lessard, A posteriori verification of invariant objects of evolution equations: periodic orbits in the Kuramoto--Sivashinsky PDE, \emph{SIAM J. Appl. Dyn. Syst.} \textbf{16} (2017) 687--728; doi:10.1137/16M1073789.

\bibitem{gs1985} M. Golubitsky and I. Stewart, Hopf bifurcation in the presence of symmetry, \emph{Arch. Rational Mech. Anal.} \textbf{87} (1985) 107--165; doi:10.1007/BF00280698.

\bibitem{gss1988} M. Golubitsky, I. Stewart and D. G. Schaeffer, \emph{Singularities and Groups in Bifurcation Theory}, Vol.~II, Applied Mathematical Sciences \textbf{69}, Springer, New York, 1988; doi:10.1007/978-1-4612-4574-2.

\bibitem{hi1997} F. C. Hoppensteadt and E. M. Izhikevich, \emph{Weakly Connected Neural Networks}, Applied Mathematical Sciences \textbf{126}, Springer, New York, 1997; doi:10.1007/978-1-4612-1828-9.

\bibitem{johansson2017} F. Johansson, Arb: efficient arbitrary-precision midpoint-radius interval arithmetic, \emph{IEEE Trans. Comput.} \textbf{66} (2017) 1281--1292; doi:10.1109/TC.2017.2690633.

\bibitem{jz2012} M. A. Johnson and K. Zumbrun, Convergence of Hill's method for nonselfadjoint operators, \emph{SIAM J. Numer. Anal.} \textbf{50} (2012) 64--78; doi:10.1137/100809349.

\bibitem{capd2021} T. Kapela, M. Mrozek, D. Wilczak and P. Zgliczy\'nski, CAPD::DynSys: a flexible C++ toolbox for rigorous numerical analysis of dynamical systems, \emph{Commun. Nonlinear Sci. Numer. Simul.} \textbf{101} (2021) 105578; doi:10.1016/j.cnsns.2020.105578.

\bibitem{kz2003} T. Kapela and P. Zgliczy\'nski, The existence of simple choreographies for the $N$-body problem: a computer-assisted proof, \emph{Nonlinearity} \textbf{16} (2003) 1899--1918; doi:10.1088/0951-7715/16/6/302.

\bibitem{kato1976} T. Kato, \emph{Perturbation Theory for Linear Operators}, 2nd edition, Grundlehren der mathematischen Wissenschaften \textbf{132}, Springer, Berlin, 1976.

\bibitem{kq2022} C. Kuehn and E. Queirolo, Computer validation of neural network dynamics: a first case study, \emph{Discrete Contin. Dyn. Syst. Ser. B} \textbf{30} (2025) 2073--2093; doi:10.3934/dcdsb.2024145; arXiv:2202.05073.

\bibitem{keb2017} P. K\"ugler, M. A. K. Bulelzai and A. H. Erhardt, Period doubling cascades of limit cycles in cardiac action potential models as precursors to chaotic early afterdepolarizations, \emph{BMC Syst. Biol.} \textbf{11} (2017) 42; doi:10.1186/s12918-017-0422-4.

\bibitem{keb2018} P. K\"ugler, A. H. Erhardt and M. A. K. Bulelzai, Early afterdepolarizations in cardiac action potentials as mixed mode oscillations due to a folded node singularity, \emph{PLoS One} \textbf{13} (2018) e0209498; doi:10.1371/journal.pone.0209498.

\bibitem{pe1994} J. E. Paullet and G. B. Ermentrout, Stable rotating waves in two-dimensional discrete active media, \emph{SIAM J. Appl. Math.} \textbf{54} (1994) 1720--1744; doi:10.1137/S0036139993250683.

\bibitem{pc1998} L. M. Pecora and T. L. Carroll, Master stability functions for synchronized coupled systems, \emph{Phys. Rev. Lett.} \textbf{80} (1998) 2109--2112; doi:10.1103/PhysRevLett.80.2109.

\bibitem{rs1998} A. M. Rucklidge and M. Silber, Bifurcations of periodic orbits with spatio-temporal symmetries, \emph{Nonlinearity} \textbf{11} (1998) 1435--1455; doi:10.1088/0951-7715/11/5/015.

\bibitem{tp2006} K. H. W. J. ten Tusscher and A. V. Panfilov, Alternans and spiral breakup in a human ventricular tissue model, \emph{Am. J. Physiol. Heart Circ. Physiol.} \textbf{291} (2006) H1088--H1100; doi:10.1152/ajpheart.00109.2006.

\end{thebibliography}

\end{document}
